% =====================================================================
%  Analisis de Optimalidad y Validacion Regulatoria de Mercados P2P
%  en Colombia — DOCUMENTO DE PROCESO
%
%  Brayan Stiven Lopez-Mendez · Universidad de Narino · 2026
%
%  Narrativa: dato -> proceso -> resultado. Cada transformacion del
%  dato tiene su propia subseccion y su propia figura de impacto.
%
%  Formato heredado de reformateo/NuevoFormato/ (article 12pt, una
%  columna, capitulos simulados con titlesec), con las adiciones que
%  exige una tesis: bibliografia IEEE, indice de tablas, nomenclatura,
%  paginas preliminares en romanos y subfiguras para los pares M1|M3.
%
%  Compilar:  pdflatex main -> bibtex main -> pdflatex main x2
% =====================================================================
\documentclass[12pt]{article}
\usepackage[utf8]{inputenc}
\usepackage[spanish]{babel}
\usepackage[margin=1in]{geometry}
\usepackage{parskip} % Elimina sangría y agrega espacio entre párrafos
\usepackage{fancyhdr}
\usepackage{graphicx}
\usepackage{xcolor}
\usepackage{hyperref}
\usepackage{enumitem}
\usepackage{listings}
\usepackage{titlesec}
\usepackage{setspace}
\usepackage{ragged2e}
\usepackage{amsmath}
\usepackage{amsfonts}
\usepackage{amssymb}
\usepackage{longtable}
\usepackage{booktabs}
\usepackage{multirow}
\usepackage{array}
\usepackage{float}
\usepackage{fancyvrb}
\usepackage{xurl}

% --- Añadidos respecto a la plantilla base ---------------------------
% Pares de paneles M1 | M3: casi toda figura del documento es doble.
\usepackage{subcaption}
% Diagramas de proceso, cadenas de flujo y esquemas de medición.
\usepackage{tikz}
\usetikzlibrary{arrows.meta, positioning, shapes.geometric, calc, fit,
                backgrounds, decorations.pathreplacing}
% Cifras con la convencion espanola: miles con punto, decimales con
% coma. Se usa numprint y NO siunitx: la version de siunitx que
% distribuye MiKTeX 24.1 exige un expl3 mas nuevo que el de la propia
% distribucion, y falla con \l_siunitx_quantity_prefix_mode_str
% indefinido. numprint es estable y cubre lo que el documento necesita.
\usepackage{numprint}
\npthousandsep{.}
\npdecimalsign{,}
% \num{4839.52} imprime <<4.839,52>>; \uni{12.0}{kWh}, <<12,0 kWh>>;
% \pct{20} imprime <<20 %>>.
\newcommand{\num}[1]{\numprint{#1}}
\newcommand{\uni}[2]{\numprint{#1}\,#2}
% OJO: \pct solo funciona en modo TEXTO. babel-spanish redefine el
% simbolo de porcentaje con una macro que inspecciona \lastskip, y
% dentro de $...$ el \, produce mu-glue, lo que aborta la compilacion
% con <<Incompatible glue units>>. Escribir  $\pm$\,\pct{0.3}  y no
% $\pm\pct{0.3}$.
\newcommand{\pct}[1]{\numprint{#1}\,\%}

% Bibliografía numérica estilo IEEE.
\usepackage[numbers,sort&compress]{natbib}

% Ruta de búsqueda de imágenes (las figuras viven en figuras/)
\graphicspath{{./}}
% Figuras en la posición exacta del texto (no flotantes)
\floatplacement{figure}{H}
\usepackage{tcolorbox}
% Marcador de figura pendiente: permite redactar antes de generarla.
\newcommand{\placeholderfig}[2][0.85\linewidth]{%
  \fbox{\parbox{#1}{\centering\vspace{1.6cm}%
  \textbf{[ FIGURA PENDIENTE ]}\\[0.4em]%
  \small #2\vspace{1.6cm}}}%
}

% --- Profundidad de numeración y tabla de contenidos ---
\setcounter{secnumdepth}{5}  % Numera hasta subparagraph (5 niveles)
\setcounter{tocdepth}{3}     % Muestra en índices hasta subsubsection (3 niveles)

% --- Redefinir "Cuadro" como "Tabla" para babel español ---
\addto\captionsspanish{%
  \def\tablename{Tabla}%
}

% --- Numeración jerárquica de figuras y tablas por sección ---
\numberwithin{figure}{section}
\numberwithin{table}{section}
\numberwithin{equation}{section}

% --- Colores personalizados ---
\definecolor{darkblue}{rgb}{0.1, 0.2, 0.4}
\definecolor{sectioncolor}{rgb}{0.15, 0.35, 0.65}
\definecolor{codegray}{rgb}{0.96, 0.96, 0.96}
\definecolor{codeblue}{rgb}{0.25, 0.41, 0.88}
\definecolor{codegreen}{rgb}{0.0, 0.5, 0.0}
\definecolor{warningcolor}{rgb}{0.8, 0.4, 0.0}
\definecolor{successcolor}{rgb}{0.0, 0.6, 0.0}
\definecolor{infocolor}{rgb}{0.0, 0.4, 0.8}
% Color de las decisiones técnicas fechadas (capítulo de depuración).
\definecolor{decisioncolor}{rgb}{0.35, 0.20, 0.50}

% --- Encabezado y pie de página ---
\pagestyle{fancy}
\fancyhead[L]{\textbf{Mercados P2P en Colombia}}
\fancyhead[R]{\textbf{Documento de proceso}}
\fancyfoot[C]{\thepage}
\renewcommand{\headrulewidth}{0.4pt}
\renewcommand{\footrulewidth}{0.4pt}

% Estilo de página para inicio de capítulo (sin encabezado)
\fancypagestyle{chapterpage}{
    \fancyhf{}
    \fancyfoot[C]{\thepage}
    \renewcommand{\headrulewidth}{0pt}
    \renewcommand{\footrulewidth}{0.4pt}
}

% Estilo de página para índices (sin encabezado ni líneas)
\fancypagestyle{indexpage}{
    \fancyhf{}
    \fancyfoot[C]{\thepage}
    \renewcommand{\headrulewidth}{0pt}
    \renewcommand{\footrulewidth}{0pt}
}

% --- Hipervínculos ---
\pdfstringdefDisableCommands{%
  \def\code#1{#1}%
  \def\file#1{#1}%
  \def\CAL#1{CAL-#1}%
}
\hypersetup{
    colorlinks=true,
    linkcolor=black,
    citecolor=black,
    urlcolor=black,
    pdftitle={Analisis de Optimalidad y Validacion Regulatoria de Mercados P2P en Colombia},
    pdfauthor={Brayan Stiven Lopez-Mendez},
    hypertexnames=false
}

% --- Estilo de títulos ---
\titleformat{\section}[display]
  {\normalfont\Large\bfseries\color{black}}
  {\vspace*{1cm}\huge Capítulo \thesection}
  {10pt}
  {\Huge\bfseries\color{black}}
  [\thispagestyle{chapterpage}]
\titlespacing*{\section}
  {0pt}{0pt}{40pt}

\titleformat{\subsection}
  {\Large\bfseries\color{black}}{\thesubsection}{1em}{}[\justifying]
\titlespacing*{\subsection}
  {0pt}{3.25ex plus 1ex minus .2ex}{1.5ex plus .2ex}

\titleformat{\subsubsection}
  {\normalsize\bfseries\color{black}}{\thesubsubsection}{1em}{}[\justifying]
\titlespacing*{\subsubsection}
  {0pt}{3ex plus 1ex minus .2ex}{1.5ex plus .2ex}

\titleformat{\paragraph}
  {\normalsize\bfseries\color{black}}{\theparagraph}{1em}{}[\justifying]
\titlespacing*{\paragraph}
  {0pt}{2.5ex plus 1ex minus .2ex}{1ex plus .2ex}

\titleformat{\subparagraph}
  {\normalsize\itshape\color{black}}{\thesubparagraph}{1em}{}[\justifying]
\titlespacing*{\subparagraph}
  {0pt}{2ex plus 1ex minus .2ex}{1ex plus .2ex}

% Los anexos imprimen «Anexo A» en vez de «Capítulo A».
\newcommand{\iniciaranexos}{%
  \clearpage
  \appendix
  \setcounter{section}{0}%
  \titleformat{\section}[display]
    {\normalfont\Large\bfseries\color{black}}
    {\vspace*{1cm}\huge Anexo \thesection}
    {10pt}
    {\Huge\bfseries\color{black}}
    [\thispagestyle{chapterpage}]%
}

% --- Configuración del listado de código ---
\lstset{
  backgroundcolor=\color{codegray},
  basicstyle=\small\ttfamily,
  keywordstyle=\color{codeblue}\bfseries,
  commentstyle=\color{codegreen},
  stringstyle=\color{codegreen},
  frame=single,
  breaklines=true,
  showstringspaces=false,
  tabsize=2,
  language=Python,
  inputencoding=utf8,
  extendedchars=true,
  literate={á}{{\'a}}1 {é}{{\'e}}1 {í}{{\'i}}1 {ó}{{\'o}}1 {ú}{{\'u}}1
           {Á}{{\'A}}1 {É}{{\'E}}1 {Í}{{\'I}}1 {Ó}{{\'O}}1 {Ú}{{\'U}}1
           {ñ}{{\~n}}1 {Ñ}{{\~N}}1
}

% --- Espaciado entre líneas ---
\onehalfspacing
\setlength{\headheight}{16pt}

% --- Comandos personalizados ---
\newcommand{\code}[1]{\textcolor{codeblue}{\texttt{\small#1}}}
\newcommand{\file}[1]{\textcolor{darkblue}{\texttt{\small#1}}}
% Referencia a una decisión técnica fechada del registro (Anexo B).
% Se llama \CAL y no \cal porque \cal ya existe en LaTeX (fuente
% caligrafica de la version 2.09) y redefinirla rompe la compilacion.
\newcommand{\CAL}[1]{\textcolor{decisioncolor}{\textbf{CAL-#1}}}
% Atribución de procedencia al pie de una figura o tabla. El documento
% no publica ninguna figura sin declarar de qué artefacto salió.
\newcommand{\fuente}[1]{\\[0.2em]{\footnotesize\itshape Fuente: #1}}
% Las dos coberturas de medición, siempre con su porcentaje.
% Parrafo interpretativo bajo una figura. Es un habito del autor en
% sus informes de 2026: tras cada figura escribe un parrafo que
% empieza por <<Nota.>>, dice lo que la figura revela y el texto no,
% y cierra con la procedencia.
\newcommand{\notafig}[1]{%
  \vspace{-0.4em}%
  {\small\noindent\textbf{Nota.} #1\par}%
  \vspace{0.5em}%
}

\newcommand{\Mone}{\mbox{M1}}
\newcommand{\Mthree}{\mbox{M3}}

% --- Cajas ------------------------------------------------------------
% Reasignadas a la narrativa de proceso de este documento:
%   detallebox  -> detalle metodológico que puede saltarse
%   trampabox   -> trampa documentada del artefacto o del dato
%   lecturabox  -> cómo se lee la figura que acaba de presentarse
%   decisionbox -> decisión técnica fechada que movió las cifras
\newtcolorbox{detallebox}{
    colback=infocolor!10,
    colframe=infocolor,
    coltitle=infocolor,
    title=\textbf{Detalle metodológico},
    fonttitle=\bfseries
}

\newtcolorbox{trampabox}{
    colback=warningcolor!10,
    colframe=warningcolor,
    coltitle=warningcolor,
    title=\textbf{Trampa documentada},
    fonttitle=\bfseries
}

\newtcolorbox{lecturabox}{
    colback=successcolor!10,
    colframe=successcolor,
    coltitle=successcolor,
    title=\textbf{Cómo leer esta figura},
    fonttitle=\bfseries
}

\newtcolorbox{decisionbox}[1]{
    colback=decisioncolor!8,
    colframe=decisioncolor,
    coltitle=decisioncolor,
    title=\textbf{Decisión #1},
    fonttitle=\bfseries
}

% =====================================================================
\begin{document}
\sloppy
\justifying

% --- PORTADA ---------------------------------------------------------
\thispagestyle{empty}
\begin{titlepage}
    \centering
    \vspace*{0.1cm}
    {\color{black}\rule{\textwidth}{2pt}\par}
    \vspace{0.6cm}
    {\large\textbf{UNIVERSIDAD DE NARIÑO}}\\[0.2em]
    {\normalsize Facultad de Ingeniería\\
     Maestría en Ingeniería Electrónica}\\[1.4cm]

    {\Large\textbf{Análisis de Optimalidad y Validación Regulatoria}}\\[0.35em]
    {\Large\textbf{de Mercados P2P en Colombia}}\\[0.9cm]

    {\large\itshape Del dato crudo al resultado: memoria detallada del proceso}\\[0.4cm]
    {\color{black}\rule{0.6\textwidth}{0.6pt}\par}

    \vfill

    {\large\textbf{Brayan Stiven López-Méndez}}\\
    \url{bsmendez24B@udenar.edu.co}\\
    Ingeniero Electrónico --- Estudiante de Maestría en Ingeniería Electrónica\\[1.0cm]

    {\normalsize\textbf{Asesores}}\\[0.3em]
    {\large\textbf{Andrés Darío Pantoja Bucheli, PhD.}}\\
    {\large\textbf{Germán Darío Obando Bravo, PhD.}}\\
    Departamento de Electrónica --- Universidad de Nariño

    \vspace{1.2cm}
    {\normalsize San Juan de Pasto, Nariño --- 2026}
    \vspace{0.6cm}
\end{titlepage}

% --- PÁGINAS PRELIMINARES (numeración romana) ------------------------
\newpage
\pagenumbering{roman}

\renewcommand{\contentsname}{\Huge\textbf{Índice general}\vspace{-0.1cm}}
\renewcommand{\listfigurename}{\Huge\textbf{Índice de figuras}\vspace{-0.1cm}}
\renewcommand{\listtablename}{\Huge\textbf{Índice de tablas}\vspace{-0.1cm}}

\thispagestyle{indexpage}
\vspace*{1cm}
\tableofcontents
\newpage

\thispagestyle{indexpage}
\vspace*{1cm}
\listoffigures
\newpage

\thispagestyle{indexpage}
\vspace*{1cm}
\listoftables
\newpage

% =====================================================================
%  Nomenclatura y acronimos — paginas preliminares, sin numerar.
%  Se completa al cierre, cuando el vocabulario del documento este fijo.
% =====================================================================
\thispagestyle{indexpage}
\vspace*{1cm}
\section*{\Huge\textbf{Nomenclatura}}
\addcontentsline{toc}{section}{Nomenclatura}
\justifying

\begin{longtable}{@{}>{\raggedright\arraybackslash}p{0.20\textwidth}>{\raggedright\arraybackslash}p{0.74\textwidth}@{}}
\toprule
\textbf{Símbolo} & \textbf{Significado} \\
\midrule
\endfirsthead
\toprule
\textbf{Símbolo} & \textbf{Significado} \\
\midrule
\endhead
\bottomrule
\endfoot
\bottomrule
\endlastfoot
M1 & Frontera de medición por el circuito principal o de inyección de
cada institución. La razón entre generación y consumo medido es del
19,1 \%. \\
\addlinespace
M3 & Frontera de medición por el circuito secundario de cada
institución. La razón entre generación y consumo medido es del 91,2 \%. \\
\addlinespace
\end{longtable}


\clearpage
\pagenumbering{arabic}
\setcounter{page}{1}
\thispagestyle{fancy}

\newcommand{\sectionbreak}{\clearpage}

% =====================================================================
%  PARTE 0 — APERTURA
% =====================================================================
% =====================================================================
%  Introducción y pregunta regulatoria
%  Parte 0 — Apertura
%
%  ENCARGO: Contexto, problema, pregunta-columna, objetivos y actividades, aportes. Cierra con el mapa del proceso completo, que prefigura el documento entero. Sin una sola cifra de resultado.
%
%  Reglas del documento:
%   - Una idea = una subseccion = una figura.
%   - Toda figura declara su procedencia con.
%   - Toda cifra sale del canon de agosto; ver GUIA.md.
%   - M1 y M3 se reportan siempre en paralelo.
%   - Ningun superlativo sin releer el CSV de la frontera correcta.
% =====================================================================
\section{Introducción y pregunta regulatoria}
\label{sec:intro}
\justifying

\placeholderfig{Capítulo pendiente de redacción. Contexto, problema, pregunta-columna, objetivos y actividades, aportes. Cierra con el mapa del proceso completo, que prefigura el documento entero. Sin una sola cifra de resultado.}


% =====================================================================
%  PARTE I — EL DATO
% =====================================================================
% =====================================================================
% El dato crudo: cinco instituciones, veintisiete fuentes
% Parte I. El dato
%
% ENCARGO: que se midio, donde y con que. 20 medidores mas 7 inversores.
% Cobertura temporal y justificacion del horizonte de 6.144 h.
%
% Reglas del documento:
% - Una idea = una subseccion = una figura.
% - Toda figura declara su procedencia con.
% - Toda cifra sale del canon de agosto; ver GUIA.md.
% - M1 y M3 se reportan siempre en paralelo.
% - Ningun superlativo sin releer el CSV de la frontera correcta.
% =====================================================================
\section{Las fuentes de medición: cinco instituciones, veintisiete equipos}
\label{sec:dato}
\justifying

Todo lo que este documento afirma descansa sobre una base de medidas
tomadas en cinco instituciones de San Juan de Pasto entre abril y
diciembre de 2025. Antes de describir qué se hizo con esas medidas hay que precisar
qué son, cuántas hay y de qué calidad, porque la solidez de
cualquier conclusión posterior está delimitada por la de este punto de
partida.

\subsection{Las cinco instituciones}
\label{sub:dato-instituciones}

La comunidad estudiada la forman cinco entidades del municipio de Pasto,
todas ellas con sistema fotovoltaico instalado y medición avanzada: la
Universidad de Nariño, la Universidad Mariana, la Universidad Cooperativa
de Colombia (campus Pasto), el Hospital Universitario Departamental de
Nariño E.S.E. y la Universidad CESMAG. A lo largo del documento se
las nombra de forma abreviada como Udenar, Mariana, UCC, HUDN y CESMAG, y
ese orden se conserva en todas las tablas y figuras.

Estas cinco son las entidades que el proyecto MTE instrumentó y de las
que, en consecuencia, existen series numéricas de generación y de consumo.
La comparación se construye alrededor de ellas.

\begin{guardado}
\begin{trampabox}
Que el conjunto venga dado tiene una consecuencia que reaparece en el
Capítulo~\ref{sec:esc} y que no conviene descubrir allí: son cinco, y
cinco es un número pequeño con efectos normativos. El artículo 20 de la
Resolución CREG 101 072 de 2025 reserva su tratamiento más favorable, el
del numeral 1, a las comunidades que cumplen tres condiciones, y la
tercera exige que el Porcentaje de Distribución de Excedentes sea
«inferior al 10\% para cada uno de los usuarios». Ese porcentaje lo
acuerdan los propios integrantes, pero el artículo 19 obliga a que sume
exactamente el 100\%, de modo que la elección no es libre: si todos los
usuarios quedaran por debajo del 10\%, la suma sería inferior a 10 veces
el número de usuarios, y alcanzar el 100\% exigiría entonces 11
participantes o más. Con cinco, la condición resulta aritméticamente
inalcanzable, y no por una característica de esta comunidad sino por su
tamaño. Tampoco hubo margen para ampliar el conjunto ni para controlar su
composición, de modo que ninguna conclusión de este trabajo se extiende
sin más a comunidades de otro tamaño o de otra mezcla de usos.
\end{trampabox}
\end{guardado}

El sistema instalado en la Universidad de Nariño ocupa el Bloque
Sur de la sede Torobajo. El de la Universidad Mariana está en el Campus Deportivo Alvernia, en el sector
occidental de la ciudad, y de ahí que los equipos que lo miden lleven ese
nombre. Los de la Universidad Cooperativa de Colombia y el HUDN sí están
en sus sedes principales, en el sector Torobajo y en el Parque Bolívar
respectivamente. El de la Universidad CESMAG es el único fuera del casco
urbano: está en el Campus Universitario San Damián, en el
corregimiento de Catambuco. La Figura~\ref{fig:dato-mapa} sitúa
las cinco sobre el territorio.

\begin{figure}[H]
    \centering
    \includegraphics[width=0.98\textwidth]{figuras/f2_01_mapa_ubicacion.pdf}
    \caption[Localización de las cinco instituciones]{Localización de las
    instituciones aliadas del proyecto MTE en el municipio de Pasto. El
    área sombreada en rosa es el perímetro urbano; el fondo verde, el
    modelo digital de elevación.
    }
    \label{fig:dato-mapa}
\end{figure}

De ese conjunto interesa una característica, que no se supone sino que se
mide: las cinco entidades no comparten rutina de consumo. Cuatro son
instituciones educativas con jornada diurna y calendario académico,
mientras que la quinta es un hospital que no cierra nunca. Esa disparidad
es la que hace que en una misma hora unas tengan excedente y otras
déficit, que es la premisa material de cualquier intercambio. El
Capítulo~\ref{sec:prep} la muestra medida, institución por institución.

\begin{guardado}
\begin{detallebox}
Son dos proximidades distintas y se confunden con facilidad. La
Figura~\ref{fig:dato-mapa} acredita la \emph{geográfica}, y es
considerable: cuatro de las cinco caben en unos pocos kilómetros. De ahí
no se sigue la \emph{eléctrica}, que es la que importaría para hablar de
un mismo circuito o de una misma subestación, y que el proyecto no
documenta. Este trabajo no la afirma. Tampoco haría falta: el modelo del
Capítulo~\ref{sec:modelo} no representa la red de distribución ni calcula
flujos de potencia, de modo que la distancia eléctrica entre participantes
no interviene en ningún resultado.

Tampoco comparten comercializador. La mayoría de las cinco tiene contrato
con un comercializador distinto del operador de red, y aun así este
trabajo liquida las cinco contra la tarifa publicada por CEDENAR. Es una
decisión tomada por disponibilidad de la fuente, y el
Capítulo~\ref{sec:precios} la declara junto con su alcance.
\end{detallebox}
\end{guardado}

Sí interviene, en cambio, una asimetría regulatoria que el
Capítulo~\ref{sec:precios} desarrolla: dos de las cinco, Udenar y el HUDN,
son entidades públicas, y las otras tres, privadas. Las dos categorías
pagan tarifas distintas por la misma energía.

\subsection{Veintisiete fuentes, dos magnitudes}
\label{sub:dato-fuentes}

Cada institución está instrumentada con cuatro medidores eléctricos y al
menos un inversor fotovoltaico, y todos registran con resolución
de dos minutos. El inventario de los equipos arroja
27 dispositivos: 20 medidores y 7 inversores. Estos últimos
están repartidos de forma desigual, porque Udenar tiene tres y las demás
uno cada una. Hay además cinco estaciones meteorológicas instaladas, una por
institución, que este trabajo no emplea.

De cada equipo se toma una sola magnitud, y conviene subrayar cuánto se
descarta al hacerlo. Un medidor registra 54
magnitudes por lectura, entre ellas la corriente y la tensión de cada fase,
la frecuencia, el factor de potencia y la distorsión armónica; de todas
ellas el modelo usa una, la potencia activa total (kW). Un inversor
registra 18, entre ellas la potencia y la tensión del lado de
corriente continua; el modelo usa una, la potencia de corriente alterna
entregada (W).

La elección merece sustento, porque el mismo medidor ofrece otras dos
potencias, la reactiva y la aparente. Lo que el mercado intercambia es
energía, y la energía es la integral de la potencia en el tiempo,

\begin{equation}
 E \;=\; \int_{t}^{t+\Delta t} p(\tau)\,\mathrm{d}\tau
   \;\simeq\; \overline{P}\,\Delta t,
 \label{eq:dato-energia}
\end{equation}

de modo que la magnitud buscada es aquella cuya integral da los kilovatios
hora que después se facturan. De las tres potencias que el equipo mide,
solo la activa lo cumple, porque es la única que transfiere energía neta a
la carga. La reactiva oscila entre la fuente y la carga, y su
transferencia a lo largo de un ciclo se anula, de modo que su integral
sobre la hora no es energía consumida. La aparente combina las dos en una
sola magnitud y por eso tampoco sirve: no permite separar la parte que sí
se consume. Quedarse con la activa no sale gratis, porque la reactiva
también se factura, y el Anexo~\ref{sec:no-representado} mide ese cargo y
muestra que no altera el ordenamiento.

Por la misma razón se toma del inversor la potencia del lado de corriente
alterna y no la del lado de corriente continua, que el equipo también
registra. Lo que la institución puede consumir o vender es lo que sale del
inversor ya convertido, después de las pérdidas de conversión, mientras
que la lectura del lado continuo describe lo que entrega el arreglo
fotovoltaico y sobrestimaría la generación disponible.

Las Tablas~\ref{tab:dato-var-medidor} y~\ref{tab:dato-var-inversor} recogen todas las variables que entrega cada clase de equipo, agrupado por
familia y con la unidad de cada magnitud.

\begingroup
\footnotesize
\begin{longtable}{@{}>{\raggedright\arraybackslash}p{9.6cm} l@{}}
\caption{Las 54 magnitudes que entrega un medidor, por
familia. La marca de tiempo no se cuenta, porque no es una magnitud
medida.}\label{tab:dato-var-medidor}\\
\toprule
\textbf{Variable} & \textbf{Unidad} \\
\midrule
\endfirsthead
\multicolumn{2}{@{}l}{\itshape Tabla \thetable, continuación}\\
\toprule
\textbf{Variable} & \textbf{Unidad} \\
\midrule
\endhead
\midrule
\multicolumn{2}{r@{}}{\itshape continúa en la página siguiente}\\
\endfoot
\bottomrule
\endlastfoot
    \multicolumn{2}{@{}l}{\textit{Estado de la red}} \\*
    Frecuencia                                        & Hz \\
    Estado de salida & --- \\
    \addlinespace
    \multicolumn{2}{@{}l}{\textit{Corriente}} \\*
    Corriente de fase (A, B, C)                       & A \\
    Corriente máxima de fase (A, B, C)                & A \\
    Corriente de neutro                               & A \\
    \addlinespace
    \multicolumn{2}{@{}l}{\textit{Tensión}} \\*
    Tensión de fase (A, B, C)                         & V \\
    \addlinespace
    \multicolumn{2}{@{}l}{\textit{Potencia activa}} \\*
    \textbf{Potencia activa total} \hfill \textbf{(la que usa el modelo)}
                                                      & kW \\
    Potencia activa de fase (A, B, C)                 & kW \\
    Demanda de potencia activa vigente                & kW \\
    Demanda máxima de potencia activa                 & kW \\
    Potencia activa acumulada                         & kW \\
    \addlinespace
    \multicolumn{2}{@{}l}{\textit{Potencia reactiva}} \\*
    Potencia reactiva total                           & kvar \\
    Potencia reactiva de fase (A, B, C)               & kvar \\
    \addlinespace
    \multicolumn{2}{@{}l}{\textit{Potencia aparente}} \\*
    Potencia aparente total                           & kVA \\
    Potencia aparente de fase (A, B, C)               & kVA \\
    Demanda de potencia aparente vigente              & kVA \\
    Demanda máxima de potencia aparente               & kVA \\
    Potencia aparente acumulada                       & kVA \\
    Registro de potencia aparente, palabras baja y alta & --- \\
    \addlinespace
    \multicolumn{2}{@{}l}{\textit{Factor de potencia}} \\*
    Factor de potencia total                          & --- \\
    Factor de potencia de fase (A, B, C)              & --- \\
    Factor de potencia en la máxima aparente importada & --- \\
    \addlinespace
    \multicolumn{2}{@{}l}{\textit{Registros de energía acumulada}} \\*
    Energía activa importada, palabras baja y alta    & kWh \\
    Energía activa exportada, palabras baja y alta    & kWh \\
    Energía reactiva positiva, palabras baja y alta   & kvarh \\
    Energía reactiva negativa, palabras baja y alta   & kvarh \\
    \addlinespace
    \multicolumn{2}{@{}l}{\textit{Calidad de la onda}} \\*
    Distorsión armónica de corriente (A, B, C)        & \% \\
    Distorsión armónica de tensión (A, B, C)          & \% \\
    Distorsión de demanda de corriente (A, B, C)      & \% \\
\end{longtable}
\endgroup

\begin{table}[H]
    \centering
    \footnotesize
    \caption{Las 18 variables que entrega un inversor, por familia.}
    \label{tab:dato-var-inversor}
    \begin{tabular}{@{}>{\raggedright\arraybackslash}p{9.6cm} l@{}}
    \toprule
    \textbf{Variable} & \textbf{Unidad} \\
    \midrule
    \multicolumn{2}{@{}l}{\textit{Lado de corriente continua}} \\*
    Potencia                                          & W \\
    Corriente                                         & A \\
    Tensión                                           & V \\
    \addlinespace
    \multicolumn{2}{@{}l}{\textit{Lado de corriente alterna, potencia}} \\*
    \textbf{Potencia activa entregada} \hfill \textbf{(la que usa el modelo)}
                                                      & W \\
    Potencia aparente                                 & VA \\
    Potencia reactiva                                 & var \\
    Factor de potencia         & \% \\
    \addlinespace
    \multicolumn{2}{@{}l}{\textit{Lado de corriente alterna, medidas eléctricas}} \\*
    Frecuencia                                        & Hz \\
    Corriente total                                   & A \\
    Corriente de fase (A, B, C)                       & A \\
    Tensión de fase a neutro (A, B, C)                & V \\
    Tensión entre fases (AB, BC, CA)                  & V \\
    \bottomrule
    \end{tabular}
\end{table}

\begin{guardado}
\begin{detallebox}
Que el modelo consuma 2 de las 72 magnitudes disponibles no
es un desperdicio sino una consecuencia de su formulación: el mercado que
se simula intercambia energía, y para eso bastan la que cada institución
consume y la que produce. Las variables restantes describen la
\emph{calidad} de esa energía, que este trabajo no modela. Quedan
disponibles para extensiones que sí la consideren, y las estaciones
meteorológicas, para trabajos que quieran predecir la generación en lugar
de tomarla medida.
\end{detallebox}
\end{guardado}

\subsubsection*{Las magnitudes de cada equipo}

Hasta aquí los equipos se han citado con el nombre que les da la plataforma,
esto es, Medidor 1 o Medidor 3, y ese nombre no dice a qué circuito se
conectaron. El inventario de instalación del proyecto sí lo dice: cruza
cada equipo con el circuito que mide, el tablero donde está montado, el
transformador de corriente y la escala de conversión. La
Tabla~\ref{tab:dato-inventario} recoge de ahí los dos medidores que el
modelo emplea en cada institución. De ese cruce sale el nombre de las dos
fronteras de medición: las nombra el circuito que miden, no el número del
equipo.

\begin{table}[H]
    \centering
    \small
    \caption{Circuito medido por los dos equipos que el modelo emplea en
    cada institución.}
    \label{tab:dato-inventario}
    \begin{tabular}{@{}l >{\raggedright\arraybackslash}p{4.3cm}
                        >{\raggedright\arraybackslash}p{4.3cm}@{}}
    \toprule
    \textbf{Institución} & \textbf{Medidor 1 (M1)} &
    \textbf{Medidor 3 (M3)} \\
    \midrule
    Udenar   & Circuito principal, tablero de la subestación del bloque Sur
             & Circuito secundario, piso 1A \\
    \addlinespace
    Mariana  & Circuito de inyección, totalizador principal del campus
               Alvernia
             & \emph{No existe.} Se aproxima con el Medidor 1 escalado \\
    \addlinespace
    UCC      & Circuito de inyección, totalizador principal del bloque A
             & Totalizador del piso 2 \\
    \addlinespace
    HUDN     & Circuito de inyección principal, subestación del piso 3
             & Alimentación ininterrumpida de ginecología \\
    \addlinespace
    CESMAG   & Circuito principal, totalizador del bloque B
             & Circuito principal, totalizador del bloque A \\
    \bottomrule
    \end{tabular}
\end{table}

De cada institución el modelo lee uno de esos dos medidores, y el trabajo
entero se recorre dos veces, una con cada elección. Los dos recorridos se
llaman en adelante M1 y M3, por el número del medidor que cada uno emplea,
y el Capítulo~\ref{sec:frontera} desarrolla en qué se diferencian y por
qué se reportan siempre los dos.

\begin{guardado}
\begin{detallebox}
No todas las fuentes instrumentadas se usan. De los 20 medidores, el
modelo lee un equipo por institución en cada frontera, de modo que 9
participan y 11 quedan como auxiliares. De los 7 inversores, 5
definen la generación que ve el modelo, uno por institución, y los otros 2
son de Udenar: intervienen en la reconstrucción de su demanda, y uno de
ellos sirve además de referencia para extender hacia atrás la serie del
inversor del proyecto, que no entró en servicio hasta septiembre. La razón de no sumar todo lo
disponible se explica en el Capítulo~\ref{sec:prep}: los equipos miden
puntos distintos del circuito, y sumarlos mezclaría totalizadores con
ramales.
\end{detallebox}

\subsection{Las magnitudes excluidas}
\label{sub:dato-reactiva}

Que la energía reactiva no entregue trabajo no significa que no se pague.
La regulación colombiana cobra el consumo reactivo inductivo que supera la
mitad de la energía activa entregada en el mismo periodo horario, y lo
factura como si fuera activa dentro de los cargos por uso de las redes.
Medido sobre el horizonte, ese exceso aparece en el \pct{15,0} de las
horas con consumo del circuito principal y en el \pct{65,7} de las del
secundario, y equivale a \uni{1865016}{COP} y \uni{1210145}{COP}. La
medición no procede de los contadores de energía reactiva del equipo, que
llegan en cero durante todo el horizonte, sino de integrar hora a hora la
potencia reactiva, del mismo modo que se hace con la activa. La
primera cifra supera el margen que separa a los dos mecanismos mejor
situados en esa frontera. Aun así el ordenamiento no cambia, porque es un
cargo que todos los mecanismos pagan sobre el mismo consumo, y el efecto
propio del mercado sobre él vale el \pct{6,8} de ese margen. El
Anexo~\ref{sec:no-representado} recoge la medición completa y las razones
por las que el trabajo se hace solo sobre energía activa.
\end{guardado}

\subsection{Cobertura temporal de las fuentes}
\label{sub:dato-gantt}

La segunda pregunta después de «qué se midió» es «desde cuándo». Los
equipos no entraron en servicio a la vez, y esa asincronía condiciona el
horizonte del estudio. La Figura~\ref{fig:dato-gantt} sitúa cada fuente
sobre el eje temporal y deja ver el escalonamiento con que la comunidad
se fue instrumentando.

\begin{figure}[H]
 \centering
 \includegraphics[width=0.95\textwidth]{figuras/f2_02_gantt_fuentes.png}
 \caption[Cobertura temporal de las 27 fuentes]{Periodo cubierto por
 cada una de las 27 fuentes instrumentadas, agrupadas por institución. La
 banda vertical es el horizonte del estudio y el anillo señala la fuente
 cuyo primer registro fija su inicio.}
 \label{fig:dato-gantt}
\end{figure}

\subsection{Delimitación del horizonte de estudio}
\label{sub:dato-horizonte}

El inicio del horizonte no se eligió: viene impuesto. Un mercado entre
cinco participantes no puede simularse mientras falte uno, de modo que la
primera fecha utilizable es aquella en que empieza a registrar la última de
las fuentes esenciales. Son esenciales 14 de las 27 fuentes instrumentadas:
el medidor de cada frontera en cada institución, nueve en total, porque el
de Mariana sirve a las dos, y la fuente que provee la generación de cada
una desde el primer día.

En cuatro instituciones esa fuente es su único inversor. En Udenar es el
inversor de referencia con el que se extiende la serie del inversor del
proyecto, que no registra hasta septiembre. En la
Figura~\ref{fig:dato-gantt} las esenciales son las barras en color, y el
anillo señala la última en registrar.

Esa última es el inversor del HUDN, que empieza a registrar el 4 de abril
de 2025 a las 05:58. Que sea un inversor y no un medidor no es un detalle
menor: dentro del HUDN los dos medidores empiezan la víspera, de modo que
lo que faltaba para completar la comunidad era la generación. Quedan fuera
de la cuenta los otros dos equipos de Udenar: el inversor del proyecto,
cuya serie se reconstruye hacia atrás y por tanto no condiciona el
arranque, y el segundo Fronius, que solo interviene en la reconstrucción
de la demanda. Contar el primero señalaría como causa del arranque a un
equipo que empieza en septiembre y que no lo causa.

De ahí esa fecha como primer día del horizonte, el primero en que las cinco
instituciones tienen registro en sus medidores y en su inversor. Se cierra
el 16 de diciembre del mismo año, lo que da \num{256} días y \num{6144}
horas.

\subsection{Completitud del registro por fuente}
\label{sub:dato-cobertura}

Queda la tercera pregunta: dentro de ese horizonte, ¿cuántas horas faltan?
La respuesta debe darse por separado para cada clase de equipo, porque
medirlas igual sería un error de método: contar la noche como dato faltante
en un inversor confundiría la ausencia de sol con la ausencia de medición.
La Figura~\ref{fig:dato-cobertura} separa las dos clases y cuenta, en cada
fuente, las horas sin registro. El denominador es el horizonte y no la
serie completa del equipo, que llega hasta 2026: se cuentan solo las horas
del horizonte en que la fuente ya estaba registrando.

\begin{figure}[H]
 \centering
 \includegraphics[width=0.95\textwidth]{figuras/f2_03_cobertura_fuentes.png}
 \caption[Horas sin registro por fuente]{Horas sin registro de cada
 fuente. Los medidores se evalúan sobre las 24 horas del día; los
 inversores, solo sobre la franja diurna, que va de las 06:00 a las 19:00.}
 \label{fig:dato-cobertura}
\end{figure}

Medidos contra las \num{6144} horas del horizonte, dos equipos de Udenar
parecerían defectuosos: el Medidor 4 marcaría el \pct{88.7} y el Inversor
MTE, el \pct{39.3} de las horas de sol. Ninguno de los dos falla. El
primero se instaló el 29 de abril y el segundo entró en septiembre, de modo
que a los dos se les contaban como huecos horas en que el equipo todavía no
existía.

Descontadas esas horas, las 27 fuentes caen entre el \pct{96.7} y el
\pct{98.8} y no queda excepción que acotar. El peor caso son las \num{204}
horas que le faltan al Medidor 4 de Cesmag, que es auxiliar; de los equipos
que el modelo emplea, a ninguno le faltan más de \num{126}. Son pocas, y
son exactamente las que el capítulo siguiente tiene que rellenar.

\subsection{Alcance y límites del inventario}
\label{sub:dato-alcance}

Antes de pasar al tratamiento es importante delimitar el alcance de esta base de
datos, porque determina lo que el resto del documento puede afirmar.

Se tienen nueve meses continuos de medición horaria sobre cinco
participantes reales, con cobertura de adquisición superior al \pct{97} en
todos los equipos que el modelo emplea y con un horizonte fijo. Son
medidas tomadas en instalaciones en operación, no series sintéticas ni
perfiles tipo.

Lo que definitivamente no se abarcó: el escenario tiene cinco
participantes, y en ningún momento este trabajo se extiende sin más a
comunidades de otro tamaño. El
Capítulo~\ref{sec:esc} muestra que ese número tiene además consecuencias
regulatorias directas, porque algunos umbrales normativos se definen sobre
la participación relativa de cada miembro y con cinco miembros ciertos
umbrales resultan matemáticamente inalcanzables. Tampoco sostiene conclusiones
estacionales completas: nueve meses cubren buena parte del ciclo anual de
Pasto, pero no un año entero.

Hay una segunda ausencia, y es de otra naturaleza: el modelo no representa
la red. No hay impedancias, ni distancias eléctricas, ni límites de
capacidad entre los nodos, de modo que las cinco instituciones entran como
cinco perfiles de carga y de generación sin más relación entre sí que la
temporal.

Esa ausencia no es un descuido sino el alcance del objeto que se simula.
Lo que la regulación colombiana define para una comunidad energética es
una liquidación en kilovatios hora sobre un periodo, no un despacho
físico: el intercambio entre pares se salda a través de la red del
operador, y ninguna de las resoluciones que el Capítulo~\ref{sec:esc}
compara exige un flujo de potencia para calcular lo que cada miembro paga.
El modelo reproduce esa liquidación y no otra cosa.

La consecuencia hay que decirla con el mismo cuidado. Los resultados
describen lo que la liquidación regulatoria arrojaría, no si la red física
podría alojar esos intercambios, y un despliegue real exigiría un estudio
de flujo de carga que este trabajo no hace. De esa familia de restricciones
el documento solo mide la que tiene precio en la factura, que es el consumo
de energía reactiva, y lo hace en el
Anexo~\ref{sec:no-representado}.

En resumen, de las 27 fuentes instrumentadas el modelo emplea 16: los
nueve medidores de frontera, los cinco inversores que definen la generación
y los dos que Udenar aporta a la reconstrucción de su demanda. Sobre esa
base se apoya el tratamiento que describe el capítulo siguiente.

% =====================================================================
% El preprocesamiento del dato
% Parte I. El dato
%
% ENCARGO: el pipeline etapa por etapa. Es el capitulo central del
% documento: cada etapa con su figura de antes y despues.
%
% Reglas del documento:
% - Una idea = una subseccion = una figura.
% - Toda figura declara su procedencia con.
% - Toda cifra sale del canon de agosto; ver GUIA.md.
% - M1 y M3 se reportan siempre en paralelo.
% - Ningun superlativo sin releer el CSV de la frontera correcta.
% =====================================================================
\section{El preprocesamiento del dato}
\label{sec:prep}
\justifying

El capítulo anterior describió qué se midió y con qué instrumentos. Este
detalla lo que le ocurre a esas lecturas antes de que el modelo las
reciba.

La distancia entre los dos extremos es amplia: de un lado están los
archivos que escribe cada uno de los 27 equipos, con una lectura cada dos
minutos y huecos allí donde la adquisición falló. Del otro lado está lo
que el modelo necesita para operar, que es cuánto consumió y cuánto generó
cada institución en cada una de las \num{6144} horas del horizonte. Entre
ambos median seis etapas, y ninguna es automática: en todas se decidió
algo. Juzgar la solidez de las cifras del
Capítulo~\ref{sec:res} exige, por tanto, haber recorrido antes ese
proceso completo.

\subsection{Estructura del canal de preprocesamiento}
\label{sub:prep-mapa}

Antes de precisar cada etapa conviene ver el flujo completo de
procesamiento, porque las seis no son todas de la misma clase: tres
alteran el valor del dato y las otras tres lo preparan, lo ensamblan o lo
verifican. La
Figura~\ref{fig:prep-pipeline} las ordena y las distingue.

\begin{figure}[H]
 \centering
 \begin{tikzpicture}[
 x=1cm, y=1cm,
 font=\linespread{1}\scriptsize\selectfont,
 etapa/.style={anchor=west, draw=black!45, fill=white,
 rounded corners=2pt, minimum width=2.30cm, minimum height=1.12cm,
 align=center, inner sep=2pt},
 clave/.style={etapa, draw=infocolor, fill=infocolor!14, thick},
 salida/.style={etapa, draw=successcolor, fill=successcolor!12,
 thick, minimum width=4.96cm},
 chip/.style={anchor=west, rounded corners=2pt, thick,
 minimum width=1.10cm, minimum height=0.42cm, inner sep=1pt},
 flecha/.style={-{Stealth[length=4.5pt, width=3.5pt]}, draw=black!55,
 line width=0.5pt},
 retorno/.style={flecha, draw=black!40, rounded corners=3pt},
 rotulo/.style={anchor=east, align=right, text=black!62,
 font=\linespread{1}\scriptsize\selectfont\itshape},
 nota/.style={align=left, text=black!62,
 font=\linespread{1}\scriptsize\selectfont\itshape},
 ]
 % Tres carriles. El del medio reúne las etapas que alteran el valor del
 % dato; los otros dos, la que lo prepara y las que lo verifican. Se
 % pintan al fondo para que las cajas queden encima.
 \begin{scope}[on background layer]
 \fill[black!4, rounded corners=3pt] (2.80,-0.72) rectangle (6.55,0.72);
 \fill[infocolor!8, rounded corners=3pt] (2.80,-2.67) rectangle (16.05,-1.23);
 \fill[black!4, rounded corners=3pt] (2.80,-4.62) rectangle (16.05,-3.18);
 \end{scope}

 % --- Carril 1: del archivo a la serie horaria ------------------------
 \node[etapa, minimum width=3.05cm] (e1) at (2.9500, 0)
 {\textbf{1}\\Del archivo a\\la serie horaria};

 % --- Carril 2: las tres etapas que alteran el valor ------------------
 \node[clave] (e2) at (2.9500, -1.95) {\textbf{2}\\Clasificar\\el medidor};
 \node[clave] (e3) at (5.6125, -1.95) {\textbf{3}\\Reconstruir la\\demanda bruta};
 \node[clave] (e4) at (8.2750, -1.95) {\textbf{4}\\Atípicos\\y huecos};

 % --- Carril 3: ensamblar, verificar y entregar -----------------------
 \node[etapa] (e5) at (2.9500, -3.90) {\textbf{5}\\Ensamblar\\y verificar};
 \node[etapa] (e6) at (5.6125, -3.90) {\textbf{6}\\Fijar la\\zona horaria};
 \node[salida] (out) at (10.9375,-3.90)
 {\textbf{Matrices $D$ y $G$}\\$5 \times \num{6144}$, no negativas};

 % --- Rótulos de carril ----------------------------------------------
 \node[rotulo] at (2.62, 0) {transporta,\\ordena\\y convierte};
 \node[rotulo, text=infocolor,
 font=\linespread{1}\scriptsize\selectfont\bfseries]
 at (2.62,-1.95) {las tres etapas\\que alteran\\el valor del dato};
 \node[rotulo] at (2.62,-3.90) {ensamblan\\y verifican};

 % --- Avance dentro de cada carril ------------------------------------
 \draw[flecha] (e2) -- (e3);
 \draw[flecha] (e3) -- (e4);
 \draw[flecha] (e5) -- (e6);
 \draw[flecha] (e6) -- (out);

 % --- Paso de un carril al siguiente ----------------------------------
 \draw[retorno] (e1.south) -- (e2.north);
 \draw[retorno] (e4.south) -- ++(0,-0.42) -| (e5.north);

 % --- La bifurcación de las dos fronteras, en la etapa 1 --------------
 \node[chip, draw=sectioncolor, fill=sectioncolor!14, text=sectioncolor]
 (pm1) at (2.95, 1.34) {\textbf{M1}};
 \node[chip, draw=warningcolor, fill=warningcolor!16, text=warningcolor]
 (pm3) at (4.15, 1.34) {\textbf{M3}};
 \draw[flecha, draw=sectioncolor, line width=0.7pt]
 (pm1.south) -- (e1.north);
 \draw[flecha, draw=warningcolor, line width=0.7pt]
 (pm3.south) -- (e1.north);
 \node[nota, anchor=west, text width=8.4cm] at (6.75, 1.34)
 {aquí se separan las dos fronteras de medición:
 cambia el medidor que se lee y, con él, su tipo};

 \node[nota, anchor=center, align=center, text=infocolor]
 at (13.31,-1.95) {aquí se juega\\la fidelidad del dato};
 \end{tikzpicture}
 \caption[El pipeline de preprocesamiento en seis etapas]{Las 6
 etapas del preprocesamiento, en el orden en que se ejecutan. La
 bifurcación entre las dos fronteras de medición ocurre en la etapa~1.}
 \label{fig:prep-pipeline}
\end{figure}

Esas tres se reparten a su vez en dos grupos: la primera lleva el archivo
hasta la serie horaria, y las dos últimas la ensamblan y la verifican. Las
tres resaltadas en la figura son las que modifican el valor. Esa separación es
la que permite decir con precisión dónde se juega la fidelidad del dato.

Del mapa importa retener una propiedad. Todo el trabajo se recorre dos
veces, una por cada frontera de medición, y lo único que las distingue es
qué medidor se lee en la primera etapa. De ahí en adelante el tratamiento es
el mismo para las dos: si algo sale distinto es porque el medidor elegido
lo provoca, no porque se le aplique una regla diferente.

Cualquier diferencia entre M1 y M3 en los resultados es, por tanto,
atribuible a qué se mide y no a cómo se procesa. Hay una excepción, y
es importante definirla ahora: Mariana no tiene un segundo medidor
que leer, de modo que en la frontera secundaria se reutiliza el primero
multiplicado por un factor que el Capítulo~\ref{sec:frontera} declara.

Las seis etapas no se reparten por igual en lo que sigue. La primera ocupa
cuatro subsecciones, porque es donde el dato cambia de paso y de forma. La
segunda y la tercera ocupan dos cada una, y en los dos casos la primera de
ellas está dedicada al fenómeno que las motiva. Las tres últimas ocupan una
cada una. A partir de ahí el capítulo deja de describir etapas: las
subsecciones finales leen el resultado y declaran qué le hacen a la cifra
las decisiones tomadas por el camino.

\subsection{Del archivo a la serie horaria}
\label{sub:prep-lectura}

La primera etapa no interviene sobre el valor medido: lo transporta, lo
ordena y lo expresa en el paso y las unidades con que el modelo trabaja.
Recibe los archivos de un equipo, con una lectura cada dos minutos, y
entrega una serie de una hora. La Figura~\ref{fig:prep-archivo-serie}
recorre ese camino sobre una hora real de la UCC, es decir, sigue las
lecturas del medidor y las del inversor desde el archivo en que llegaron
hasta el valor horario que cada una deja en la serie.

\begin{figure}[H]
 \centering
 \includegraphics[width=0.95\textwidth]{figuras/f3_01_archivo_a_serie_m1.png}
 \caption[Del archivo a la serie horaria, sobre una hora real]{El
 recorrido de una hora real, la de las 13 horas del jueves 6 de noviembre
 de 2025 en la UCC, frontera M1.}
 \label{fig:prep-archivo-serie}
\end{figure}

Los dos equipos no hablan en la misma unidad. El medidor registra en
kilovatios y el inversor en vatios, de modo que sus series no se pueden
comparar ni sumar tal como llegan. La del inversor se divide entre mil y
la del medidor se usa como viene.

La etapa lleva además cuatro salvaguardas contra cosas que este
conjunto no contiene, y de ellas lo que importa es cuántas veces actúan:
ninguna.

\begin{enumerate}[leftmargin=*, itemsep=0.35em]
 \item \textbf{Archivos que no se dejan leer.} No hay ninguno entre los 73
 del árbol.

 \item \textbf{Fechas ilegibles.} No hay ninguna entre los \num{2813305}
 registros de los diez medidores.

 \item \textbf{Valores no numéricos.} Tampoco hay ninguno entre esos mismos
 registros.

 \item \textbf{Lecturas negativas del inversor.} No hay ninguna entre las
 \num{793981} de los siete equipos.
\end{enumerate}

La última no actúa porque un inversor entrega energía a la red y no la toma
de ella. Los siete registran el cero miles de veces, de modo que la
ausencia de valores negativos no se debe a que el equipo no llegue hasta
ahí, sino a que por debajo del cero no hay nada que medir.

Queda una sola ambigüedad que la etapa sí tiene que resolver, y es la del
instante que llega dos veces.

\subsection{Los instantes duplicados}
\label{sub:prep-duplicados}

Antes de promediar hay que resolver una ambigüedad del dato crudo. En la totalidad de la base de datos algunas
lecturas comparten instante con otra, \num{3108} veces en la frontera
principal y \num{3588} en la secundaria, pero el registro de un equipo no
puede llevar dos valores en el mismo instante. Esto se resuelve dejando uno solo,
tomando el promedio entre sí de los datos duplicados. En el horizonte de estudio solo quedan 27 en cada frontera.

El fenómeno es de un equipo y no del conjunto: \num{3073} de los
\num{3108} instantes de la frontera principal son del medidor de la UCC, y
en las otras cuatro instituciones no pasan de 30 apariciones en todo el
archivo. Por eso el ejemplo de la Figura~\ref{fig:prep-archivo-serie} es de
la UCC, que es la única institución donde la fusión llega a ocurrir dentro
del horizonte.

La regla para determinar qué instante de los datos duplicados usar no tiene consecuencia. En los \num{6696} instantes repetidos de las dos fronteras las
lecturas que lo comparten coinciden hasta el último dígito, de modo que
promediarlas, quedarse con la primera o quedarse con la última dan
exactamente el mismo número. Se toma el promedio porque el instante ha de ser único para poder alinear los tramos
de un mismo medidor, no porque haya dos valores en disputa que arbitrar.

La regla vale igual dentro de un archivo y entre archivos. Los tramos de un
mismo medidor comparten un instante en su costura, una vez por institución
y por frontera, y también ahí la lectura repetida se funde en una sola.

\begin{figure}[H]
 \centering
 \includegraphics[width=0.95\textwidth]{figuras/f3_01d_duplicados.png}
 \caption[Cuándo y en qué equipo se repiten los instantes]{Los instantes
 repetidos de las dos fronteras de medición, en el tiempo y por
 institución.}
 \label{fig:prep-duplicados}
\end{figure}

\subsection{La energía de la hora}
\label{sub:prep-energia}

La serie se remuestrea entonces a una hora tomando la media de sus
muestras, y esa media no es una elección de método: sale de la definición
de energía. La energía de la hora es la integral de la potencia sobre ella,
según la ecuación~\eqref{eq:dato-energia}, y de esa integral solo se
conocen valores separados por un paso fijo, de modo que el tramo observado
se suma por rectángulos de ese ancho:

\begin{equation}
 E \;=\; \int_{t}^{t+\Delta t} p(\tau)\,\mathrm{d}\tau
   \;\simeq\; \delta \sum_{k=1}^{n} P_{k}
   \;=\; \overline{P}\,(n\,\delta) ,
 \qquad \delta = \uni{2}{min},
 \label{eq:prep-media}
\end{equation}

donde $n$ es el número de muestras que la hora trae, $\delta$ el paso de
muestreo y $\overline{P}$ la media. Que el paso sea fijo no es un supuesto
sino un hecho medido: sobre \num{1117638} registros de cuatro medidores
vale exactamente dos minutos en más del
\pct{99.6} de los intervalos. De ahí que sumar los tramos y promediar las
muestras no sean dos métodos entre los que haya que elegir, sino la misma
operación escrita de dos maneras.

Con la hora completa $n$ vale 30, el producto $n\,\delta$ vale una hora, y
la energía en kilovatios hora resulta numéricamente igual a la potencia
media en kilovatios, de modo que la misma serie puede leerse de las dos
maneras y a lo largo del documento se usan ambas según convenga. La
Figura~\ref{fig:prep-energia} lo enseña sobre la hora del ejemplo.

\begin{figure}[H]
 \centering
 \includegraphics[width=0.95\textwidth]{figuras/f3_01b_energia_hora.png}
 \caption[La energía de una hora, contada de dos maneras]{Las treinta
 lecturas de la hora del ejemplo dibujadas como rectángulos de dos minutos,
 y encima el rectángulo único de la hora entera a la altura de la media. Las
 dos áreas son la misma área y los dos números coinciden hasta el último
 dígito, porque el paso de muestreo es constante.}
 \label{fig:prep-energia}
\end{figure}

\subsection{La hora incompleta}
\label{sub:prep-incompleta}

Cuando la hora llega incompleta, en cambio, los rectángulos no cubren el
intervalo y el tramo sin muestrear no queda determinado por el dato. Ese es
el único punto de la etapa en el que hay que suponer algo, y aquí se le
atribuye al tramo la potencia media de las muestras presentes, es decir, se
toma $E \simeq \overline{P}\,\Delta t$ con $\Delta t = \uni{1}{h}$. Es un
supuesto y no una identidad, de modo que se declara como tal y se acota en
la subsección~\ref{sub:prep-sesgo}.

Una hora puede llegar de tres maneras, y el tratamiento cambia en cada una.

\begin{enumerate}[leftmargin=*, itemsep=0.35em]
 \item \textbf{Completa}, con sus 30 muestras. La energía queda
 determinada por el dato y no hay nada que decidir.

 \item \textbf{Incompleta}, con entre 1 y 29 muestras. Se estima con la
 media de las presentes, que es el supuesto que se acaba de declarar. Son
 \num{835} horas del horizonte.

 \item \textbf{Vacía}, sin ninguna muestra. No se estima aquí: pasa a la
 etapa de limpieza, que interpola sus primeras horas y repite un valor en
 las restantes, según describe la subsección~\ref{sub:prep-huecos}. Son
 \num{567} horas.
\end{enumerate}

La media horaria ignora las muestras que faltan, de modo que la
estimación es lo que la agregación hace por defecto. La otra opción es contar solo los minutos medidos y no añadir nada por los
que faltan. Parece la opción prudente, la que no supone nada, pero supone
tanto como la anterior: da por hecho que en los minutos sin registro no se
consumió energía, y en instalaciones con carga de base permanente eso es
falso. La Figura~\ref{fig:prep-incompleta} pone las dos suposiciones frente
al contador de energía del propio medidor, que es el único dato capaz de
arbitrar entre ellas.

\begin{figure}[H]
 \centering
 \includegraphics[width=0.95\textwidth]{figuras/f3_01c_hora_incompleta.png}
 \caption[La hora incompleta y sus dos suposiciones]{A la izquierda, una
 hora que trae 18 de sus 30 muestras: suponer la media para el tramo que
 falta deja la hora en \uni{8.98}{kWh} y el contador del equipo marcó
 \uni{9.00}{kWh}, mientras que quedarse con lo observado la dejaría en
 \uni{5.39}{kWh}. A la derecha, la misma comparación sobre todas las horas
 incompletas de los tres medidores cuyo contador está bien escalado.}
 \label{fig:prep-incompleta}
\end{figure}

El regulador, ante el mismo problema, tampoco asigna cero. La Resolución
CREG 038 de 2014 ordena estimar la lectura mientras el sistema de medición
está en falla, y en las fronteras que no reportan al administrador del
mercado se estima por consumos promedios del mismo usuario, que es el
criterio que aquí se aplica. Rige sobre fronteras comerciales y no sobre
instrumentación de investigación como esta, de modo que se toma como
criterio y no como parte de la alineación con la normativa. El articulado y su alcance exacto quedan en el Anexo~\ref{anx:repro}.que recoge el procedimiento completo de la etapa, con esas cuatro comprobaciones y el orden exacto en que se aplican.

\begin{guardado}
\begin{detallebox}

Las dos decisiones se pueden contrastar, porque el medidor lleva junto a la
potencia sus propios contadores de energía importada y exportada, en
kilovatios hora, que avanzan por su cuenta y no intervienen en el
preprocesamiento. El contraste usa la diferencia entre los dos y no la
importada sola, por la razón que da la nota de la
Figura~\ref{fig:prep-incompleta}. Sobre las \num{5810} horas completas del
circuito principal de la UCC en las que el contador tiene lectura en los dos
bordes, lo que avanza y la media horaria difieren en el \pct{0.04}
acumulado.

El supuesto sobre el tramo no observado admite una prueba más directa, y es
la que decide. En los \num{718} cortes de telemetría de los tres medidores
cuyo contador está bien escalado, el contador avanza \uni{809}{kWh};
suponer que el equipo siguió consumiendo predice \uni{812}{kWh} y suponer
que no hubo energía predice \uni{288}{kWh}. El equipo sigue integrando
durante los cortes, de modo que atribuirles energía nula falla por física y
no por estadística, y falla cada vez peor cuanto más largo es el corte.

Hay además dos cantidades que se parecen y difieren por un factor de ocho. Las \num{835} horas incompletas del horizonte
\emph{contienen} \uni{8100}{kWh}, el \pct{3.2} de la demanda, pero de esa
energía casi toda está observada: lo que el supuesto \emph{imputa} es solo
el tramo que falta, y son \uni{1054}{kWh}, el \pct{0.41}. La razón es que
esas horas pierden muy poco. El \pct{55.6} pierde una sola muestra de las
treinta y el \pct{77} pierde cuatro o menos.

Queda por saber si la estimación descansa sobre las horas peor cubiertas, y
no lo hace. Tomando como válida solo la hora que trae al menos el \pct{75}
de sus muestras, \num{147} horas del horizonte quedarían fuera, es decir el
\pct{0.49} de las que tienen dato, y tratarlas como huecos en vez de
estimarlas mueve la demanda de la comunidad en \uni{65.6}{kWh}, el
\pct{0.026}. Ninguna institución se mueve más del \pct{0.42}, y esa es
Udenar.
\end{detallebox}
\end{guardado}



\subsection{La demanda negativa}
\label{sub:prep-negativa}

La primera etapa, la que lleva del archivo a la serie horaria, entrega
datos que no concuerdan: la demanda de tres de las cinco instituciones baja
de cero durante varias horas al día. Un edificio no consume energía
negativa. Y el tratamiento no lo causó, porque esa etapa no corrige, no
sustituye ni descarta ninguna medida, de modo que lo que la serie trae es
lo que el medidor entregó. La Figura~\ref{fig:prep-negativa} recoge el
fenómeno en sus dos dimensiones y en las dos fronteras: la forma que toma
en la institución más afectada de cada una y la hora del día en que
aparece.

\begin{figure}[H]
 \centering
 \includegraphics[width=0.95\textwidth]{figuras/f3_02_demanda_negativa.png}
 \caption[La demanda que el medidor entrega en negativo]{Una fila por
 frontera de medición. A la izquierda, el perfil de la institución que
 carga con el fenómeno en cada una, en promedio y en mínimo de cada hora,
 sobre las otras cuatro en promedio y en gris; la banda sombreada es la
 zona imposible. A la derecha, cuántas horas caen bajo cero en cada hora
 del día, apiladas por institución.}
 \label{fig:prep-negativa}
\end{figure}

\notafig{Las dos curvas de la institución destacada van juntas a
propósito. Bajo M1 el promedio horario de Udenar entra en la zona
imposible y baja a \uni{-11.0}{kW} a las 13, de modo que el neteo no
arrastra unos pocos casos sueltos sino la media entera; bajo M3 ningún
promedio la alcanza y solo el mínimo la roza. Los conteos por institución,
la fracción, el mínimo y la energía hacia la red los da la
Tabla~\ref{tab:prep-tipos}. Elaboración propia desde el caché de estados
intermedios.}

El reparto es muy desigual. En Udenar la demanda cae bajo cero en cerca de
una cuarta parte del horizonte, es decir, el \pct{24.7} de las \num{6144}
horas y el \pct{25.1} de las que traen dato. En Mariana y en la UCC el
fenómeno es ocasional, y en el HUDN y en CESMAG no aparece ni una hora, de
modo que una sola institución concentra el \pct{83} de todas las horas
afectadas. La Tabla~\ref{tab:prep-tipos} reúne el conteo y el mínimo de
cada una.

Lo decisivo, sin embargo, no es cuántas horas son sino \emph{a qué hora del
día} ocurren. Si las lecturas negativas fueran ruido del sensor se
repartirían de manera aproximadamente uniforme a lo largo del día. Lo que
se cuenta hora a hora es lo contrario: las negativas dibujan la campana de
la producción solar, con su máximo de \num{288} horas a las 12 y \num{282}
a las 13, y no hay ni una sola entre las 18 y las 6.

Ese patrón identifica la causa sin necesidad de abrir el equipo. No es un
sensor averiado, es un medidor que está restando la generación antes de
reportar, y dónde esté puesto respecto de la inyección es lo que decide si
lo hace.

Y decide también en qué frontera aparece, que es lo que separa las dos
filas de la figura. En la secundaria el fenómeno casi desaparece: de las
cinco instituciones solo Mariana registra lecturas negativas, y son las
mismas \num{213} horas de la principal, porque en M3 se lee su mismo
medidor multiplicado por un factor. Los otros cuatro circuitos secundarios
no bajan de cero ni una hora, de modo que la lectura imposible es un
asunto de la frontera principal.

\subsection{Clasificación del medidor: tres tipos}
\label{sub:prep-tipos}

La explicación de ese patrón es de topología eléctrica, o dicho de otro
modo, de en qué punto del circuito se conecta cada inversor. Cuando el
inversor entra \emph{entre el medidor y la carga}, la energía que produce
no llega a pasar por el medidor, de manera que este no registra la demanda
del circuito sino lo que queda de ella al descontarle la generación, es
decir, el flujo neto $D - G$. En las horas en que la generación supera la
carga del circuito, ese flujo se invierte de sentido y el medidor registra
un valor negativo. La lectura no es un
error: es la evidencia de un autoconsumo que el medidor no puede separar.

Según dónde entre el inversor, el tratamiento se bifurca en dos: a unos
medidores hay que devolverles la generación que descontaron y a otros no
hay nada que devolverles. El documento nombra además un caso intermedio,
el neto parcial, pero no porque reciba un tratamiento propio, que es el
mismo del neto, sino porque el neteo apenas llega a invertir el flujo. Lo
que separa a esos dos es el impacto.

\begin{figure}[H]
 \centering
 \begin{tikzpicture}[
 font=\scriptsize,
 eq/.style={draw=black!55, fill=black!3, rounded corners=1.5pt,
 minimum width=1.5cm, minimum height=0.58cm, align=center},
 med/.style={draw=warningcolor, fill=warningcolor!14, very thick,
 circle, minimum size=0.72cm, inner sep=0pt},
 inv/.style={draw=successcolor, fill=successcolor!14, very thick,
 rounded corners=1.5pt, minimum width=1.15cm,
 minimum height=0.55cm, align=center},
 lin/.style={-{Stealth[length=3.5pt]}, draw=black!60},
 et/.style={font=\scriptsize\itshape, text=black!65},
 ]
 % --- net ---
 \node[et] at (-2.35,0.62) {\textbf{neto}};
 \node[eq] (r1) at (-1.0,0) {Red};
 \node[med] (m1) at (0.55,0) {M};
 \node[eq] (c1) at (2.35,0) {Carga};
 \node[inv] (i1) at (2.35,-1.0) {Inversor};
 \draw[lin] (r1) -- (m1);
 \draw[lin] (m1) -- (c1);
 \draw[lin] (i1) -- (c1);
 \node[et, text width=3.4cm, right=0.30cm of c1, align=left]
 {el inversor entra entre\\el medidor y la carga:\\el medidor ve $D-G$};

 % --- gross ---
 \node[et] at (-2.35,-2.38) {\textbf{bruto}};
 \node[eq] (r2) at (-1.0,-3.0) {Red};
 \node[inv] (i2) at (0.55,-4.0) {Inversor};
 \node[med] (m2) at (2.05,-3.0) {M};
 \node[eq] (c2) at (3.6,-3.0) {Carga};
 \draw[lin] (r2) -- (m2);
 \draw[lin] (m2) -- (c2);
 \draw[lin] (i2) -- (r2);
 \node[et, text width=3.4cm, right=0.30cm of c2, align=left]
 {el inversor entra antes\\del medidor, que ve $D$\\sin descontar nada};
 \end{tikzpicture}
 \caption[Por qué unos medidores netean y otros no]{Posición relativa
 del medidor y del inversor en los dos casos extremos. Arriba, el inversor
 entra entre el medidor y la carga, de modo que el medidor observa el
 flujo neto, que se invierte en horas de alta irradiancia. Abajo, entra
 del lado de la red y el medidor registra la demanda bruta. El
 caso intermedio, el neto parcial, recibe el mismo tratamiento que el
 neto y se distingue solo por el grado, porque la generación que queda
 detrás del medidor pesa poco frente a la carga del circuito.}
 \label{fig:prep-tipos}
\end{figure}

\begin{table}[H]
 \centering
 \caption{Clasificación de los medidores de demanda y huella del neteo, en
 las dos fronteras. Las cuatro columnas de la derecha se leen directamente
 sobre la serie que cada medidor entrega, sin reconstruir nada: son horas
 en que su lectura fue negativa, la fracción que representan sobre las
 horas con dato de ese equipo, su lectura más baja del horizonte y la
 energía que en esas horas salió del circuito hacia la red.}
 \label{tab:prep-tipos}
 \small
 \begin{tabular}{@{}l l r r r r@{}}
 \toprule
 \textbf{Institución} & \textbf{Tipo} & \textbf{Horas} &
 \textbf{Sobre sus} & \textbf{Lectura} & \textbf{Energía hacia} \\
 & & \textbf{bajo cero} & \textbf{horas} & \textbf{mínima (kW)} &
 \textbf{la red (kWh)} \\
 \midrule
 \multicolumn{6}{@{}l}{\itshape Frontera M1} \\
 Udenar & neto & \num{1517} & \pct{25.1} & \num{-33,57} & \num{15348} \\
 Mariana & neto parcial & \num{213} & \pct{3.5} & \num{-2,41} & \num{156} \\
 UCC & neto parcial & \num{94} & \pct{1.6} & \num{-5,91} & \num{200} \\
 HUDN & bruto & \num{0} & \pct{0} & \num{6,12} & \num{0} \\
 CESMAG & bruto & \num{0} & \pct{0} & \num{0,20} & \num{0} \\
 \addlinespace
 \multicolumn{6}{@{}l}{\itshape Frontera M3} \\
 Udenar & bruto & \num{0} & \pct{0} & \num{0,22} & \num{0} \\
 Mariana & bruto & \num{213} & \pct{3.5} & \num{-0,72} & \num{47} \\
 UCC & bruto & \num{0} & \pct{0} & \num{0,73} & \num{0} \\
 HUDN & bruto & \num{0} & \pct{0} & \num{1,29} & \num{0} \\
 CESMAG & bruto & \num{0} & \pct{0} & \num{0,20} & \num{0} \\
 \bottomrule
 \end{tabular}
\end{table}

El bloque de M3 dice algo que no se deduce del de M1: en el circuito
secundario los cinco medidores son brutos, de modo que la frontera no
cambia solo cuánta carga se mide sino la semántica del propio equipo. La
única excepción aparente es Mariana, y no es tal: esa institución no tiene
medidor en el circuito secundario, de modo que se reutiliza el del
principal multiplicado por un factor, y con él hereda sus \num{213} horas
negativas, ya escaladas.

Los tres tipos no nombran tres aparatos distintos sino tres situaciones
del mismo, y lo que las separa es dónde se conecta la generación. Si entra
entre el medidor y la carga, el medidor no la ve llegar y registra solo la
demanda que queda después de restarla, de modo que su lectura puede caer
bajo cero. Si entra del lado de la red, el medidor la ve pasar y su lectura
nunca se invierte. La Figura~\ref{fig:prep-tipos} dibuja los dos casos.

La tabla ordena las cinco instituciones por lo que esa posición produce, y
la gradación se lee sin más. En Udenar el flujo se invierte una de cada
cuatro horas y llega a \uni{-33.57}{kW}. En Mariana y en la UCC ocurre en
el \pct{3.5} y el \pct{1.6} de sus horas, con mínimos veinte veces menos
profundos, de modo que la inversión es allí ocasional. En el HUDN y en
CESMAG no ocurre nunca, y la lectura más baja de CESMAG, \uni{0.20}{kW},
dice que su circuito llega a rozar el cero sin cruzarlo.

La columna de la derecha es la que mide el fenómeno sin ambigüedad, porque
no depende de ninguna reconstrucción: en el horizonte, del circuito de
Udenar salieron hacia la red \uni{15348}{kWh}, contra los \uni{156}{} de
Mariana y los \uni{200}{kWh} de la UCC.

En el HUDN y en CESMAG el inversor entra fuera del tramo que el medidor
vigila, o alimenta un circuito distinto del medido, y por eso el flujo
nunca se invierte allí aunque las dos instituciones tengan su equipo. La
Figura~\ref{fig:prep-profundidad} pone esa gradación en una sola magnitud,
es decir, sitúa a cada medidor entre la lectura que entrega de ordinario y
la más baja que llegó a entregar, y lo hace en las dos fronteras.

\begin{figure}[H]
 \centering
 \includegraphics[width=0.95\textwidth]{figuras/f3_02b_profundidad.png}
 \caption[Hasta dónde baja la lectura de cada medidor]{Una fila por
 institución en el orden fijo y un panel por frontera, con los kilovatios
 en el eje horizontal. Cada fila lleva la lectura mediana del medidor, el
 recorrido desde esa mediana hasta la más baja, y la mitad central de sus
 lecturas negativas cuando las tiene.}
 \label{fig:prep-profundidad}
\end{figure}

En Udenar esa mitad central cae entre
\uni{-15.1}{} y \uni{-3.9}{kW}, con mediana de \uni{-8.7}{kW}, de modo que
el mínimo de \uni{-33.57}{kW} casi cuadruplica la inversión ordinaria; en
Mariana la mediana de las negativas es \uni{-0.59}{kW} y en la UCC,
\uni{-1.94}{kW}. Cuántas horas resume cada fila lo da la
Tabla~\ref{tab:prep-tipos}.

En la frontera M3 la clasificación no hace falta, porque allí los cinco
medidores son brutos y ninguna institución requiere reconstrucción. La
excepción es Mariana, y no porque su circuito secundario netee: es que no
tiene un segundo medidor, de modo que su serie de M3 se construye
multiplicando la del primero por un factor, conforme declara el
Capítulo~\ref{sec:frontera}. Con la serie viajan sus \num{213} horas
negativas, reducidas por ese mismo factor, y llevarlas a cero cuesta
\uni{46.7}{kWh}, el \pct{0.33} de su demanda en esa frontera.

\subsection{La reconstrucción de la demanda bruta}
\label{sub:prep-reconstruccion}

Identificada la causa, la corrección es sencilla de enunciar: si el medidor
restó la generación, se le devuelve. El modelo necesita la demanda bruta
$D$ (lo que el circuito medido consume) y la generación $G$ por separado, de
modo que para los medidores netos y los netos parciales la demanda se
reconstruye como

\begin{equation}
 D_{\text{recon}}
 = \Bigl(D_{\text{neto}} + \sum_{\text{inversores}} G \Bigr)^{+},
 \label{eq:prep-reconstruccion}
\end{equation}

donde la suma recorre \emph{todos} los inversores físicos que inyectan
entre ese medidor y su carga, y el exponente $+$ significa que, si el
resultado sale negativo, se toma cero.

La expresión no se aplica por igual a las cinco instituciones. La
Tabla~\ref{tab:prep-reconstruccion} reparte los cinco casos: a quién se le
devuelve generación, con cuántos inversores y si el recorte a cero llega a
actuar.

\begin{table}[H]
 \centering
 \caption{Cómo se aplica la reconstrucción en cada institución. El HUDN y
 CESMAG entregan lectura bruta, de modo que no hay generación que
 devolverles y el recorte queda solo como resguardo. Cuando una hora no
 trae lectura de medidor pero sí de inversor, la demanda de partida se
 toma como nula y el resultado queda igual a la generación devuelta; pasa
 en \num{7} horas del horizonte, todas de Udenar, y suma \uni{34.1}{kWh}.}
 \label{tab:prep-reconstruccion}
 \small
 \begin{tabular}{@{}l l r r r@{}}
 \toprule
 \textbf{Institución} & \textbf{Tipo} & \textbf{Inversores} &
 \textbf{Horas con} & \textbf{Energía} \\
 & & \textbf{devueltos} & \textbf{recorte} &
 \textbf{recortada (kWh)} \\
 \midrule
 Udenar & neto & \num{3} & \num{353} & \num{1013,3} \\
 \addlinespace
 Mariana & neto parcial & \num{1} & \num{6} & \num{0,6} \\
 \addlinespace
 UCC & neto parcial & \num{1} & \num{0} & \num{0} \\
 \addlinespace
 HUDN & bruto & --- & \num{0} & \num{0} \\
 \addlinespace
 CESMAG & bruto & --- & \num{0} & \num{0} \\
 \bottomrule
 \end{tabular}
\end{table}

La única fila que pide explicación es la primera, porque los tres equipos
de Udenar no cumplen el mismo papel. El inversor del proyecto entró en
servicio el 3 de septiembre de 2025 y cubre \pct{22.9} del horizonte; los
otros dos registran desde el primer día. Antes de esa fecha, entonces, la
suma solo puede devolver lo que midieron esos dos, y queda incompleta.

Que sumar los tres sea lo correcto lo dice el propio medidor. Desde el 3
de septiembre, devolverle la generación de los tres deja
\emph{exactamente} cero horas con demanda negativa. Las \num{353} de la
tabla caen todas antes, que es donde la suma va incompleta.

Tomar como nula la demanda de partida tiene además una consecuencia de
mucho más alcance, y es que convierte en cero toda hora sin lectura de
medidor. Bajo M1 eso alcanza a \num{330} horas: \num{90} de Udenar,
\num{114} de Mariana y \num{126} de la UCC. Entran al modelo como consumo
nulo sin pasar por la etapa de limpieza, que ya no las ve, y sin que
ninguna marca las distinga de una hora realmente medida.

La Figura~\ref{fig:prep-reconstruccion} confronta las tres series sobre un
mismo eje y muestra el resultado de aplicar esa devolución. La coincidencia
temporal es el argumento: la lectura se hunde exactamente cuando la banda
de generación se ensancha, y ninguna avería de sensor produciría eso.

La Figura~\ref{fig:prep-mosaico} repite esa lectura en las diez series, de
modo que se vea a quién se le aplica la devolución y a quién no.

\begin{figure}[H]
 \centering
 \includegraphics[width=0.95\textwidth]{figuras/f3_03b_reconstruccion_mosaico.png}
 \caption[La reconstrucción en las diez series]{Perfil medio horario de la
 lectura del medidor y de la demanda reconstruida en las cinco
 instituciones y las dos fronteras, con la generación devuelta como banda
 entre las dos curvas. Las escalas verticales son independientes por
 panel.}
 \label{fig:prep-mosaico}
\end{figure}

\notafig{Siete de los diez paneles salen sin banda porque su medidor no
netea y no hay generación que devolverles, de modo que la fila de la
frontera secundaria funciona como control de la principal. La excepción es
la Universidad Mariana bajo M3, donde tampoco hay devolución y sin embargo
el recorte a cero todavía quita \uni{46.7}{kWh}, que son los
\uni{156}{kWh} de la frontera principal multiplicados por el mismo factor
con que se construye su serie secundaria. Elaboración propia desde el
caché de estados intermedios.}

\begin{figure}[H]
 \centering
 \includegraphics[width=0.95\textwidth]{figuras/f3_03_reconstruccion_m1.png}
 \caption[La reconstrucción de la demanda bruta de Udenar]{La
 reconstrucción sobre Udenar. A la izquierda, un viernes hora a hora: en
 rojo la lectura del medidor y en azul la demanda que resulta al
 devolverle la generación.}
 \label{fig:prep-reconstruccion}
\end{figure}

La banda beige entre las dos curvas es esa
generación, de modo que la curva azul es la roja levantada por la banda,
y la cota vertical la mide en la hora de lectura más baja, cuando los
tres inversores entregaron \uni{35.1}{kW}. El día pertenece al tramo
posterior al 3 de septiembre de \num{2025}, cuando los tres registran a
la vez, y ninguna de sus horas llega al recorte. A la derecha, el
promedio de las \num{6144} horas en la misma escala vertical. Abajo, la
misma igualdad en energía sobre el horizonte completo.

\begin{guardado}
\notafig{La curva azul resultante recupera la forma que cabe esperar de un
campus universitario, es decir, base nocturna estable, ascenso matinal,
valle de almuerzo y tarde sostenida. La franja inferior separa las dos
partes de la cuenta: de los \uni{33705}{kWh} en que la reconstrucción
levanta la serie de Udenar, \uni{32691}{kWh} son generación devuelta y
\uni{1013}{kWh} proceden del recorte a cero. Elaboración propia aplicando
la ecuación~\eqref{eq:prep-reconstruccion}.}


\begin{detallebox}
Los inversores intervienen en dos capas distintas. La generación $G$ que
el modelo negocia proviene de \emph{un} inversor designado por
institución, que en Udenar es el del proyecto. La reconstrucción de la
ecuación~\eqref{eq:prep-reconstruccion} suma, en cambio, \emph{todos} los
que inyectan detrás del medidor, que en Udenar son tres. Las dos capas
responden a preguntas distintas: cuál generación se negocia, y cuánta
restó el medidor.
\end{detallebox}
\end{guardado}

\begin{guardado}
\subsection{Lo que cuesta reconstruir}
\label{sub:prep-costo}

El recorte a cero de la ecuación~\eqref{eq:prep-reconstruccion} no es
gratuito, y lo que cuesta se puede medir. Actúa cuando la suma de
inversores se queda \emph{corta} frente a la generación que el medidor
descontó, porque solo entonces $D_{\text{neto}} + \sum G$ sigue siendo
negativo; si la suma sobrara, el resultado saldría más positivo, no menos.
La causa es de cobertura: antes del 3 de septiembre de \num{2025} el
inversor del proyecto no registraba, y la suma devuelve al medidor menos
de lo que este había restado.

Lo que eso elimina de la serie de demanda queda acotado. Son las
\num{353} horas de Udenar de la tabla anterior, el \pct{5.7} del
horizonte, todas en la ventana solar y con máximo a las 13 horas, que
suman \uni{1013.3}{kWh}, es decir, el \pct{2.29} de la demanda de Udenar y
el \pct{0.32} de la comunitaria. La distorsión tiene además signo
conocido, porque subestima la demanda de Udenar justo en las horas de
sol.

\begin{lecturabox}
Que el signo sea conocido importa más que la magnitud. Subestimar la
demanda de Udenar en horas solares reduce su déficit y, con él, la energía
que Udenar podría comprar en esas horas. El efecto sobre el mercado entre
pares es, por tanto, restrictivo: hace el mercado más pequeño, no más
grande. Es la primera de una serie de decisiones que apuntan todas en la
misma dirección, y que la subsección~\ref{sub:prep-sesgo} reúne.
\end{lecturabox}
\end{guardado}

\subsection{El umbral de los atípicos}
\label{sub:prep-umbral}

La última etapa que altera el valor del dato es la limpieza, que primero
retira los valores imposibles y después rellena lo que falta. Actúa sobre
las 10 series de medida, que son la demanda y la generación de cada
institución, aunque en la práctica solo sobre la demanda: las cinco de
generación atraviesan la etapa sin que se les toque una sola hora.

Todo se hace a traves del siguiente criterio, un valor se retira cuando supera el umbral dado por:

\begin{equation}
 u = \max\bigl(Q_{75} + 5\,\mathrm{IQR},\; 1{,}2 \cdot P_{99{,}5}\bigr),
 \label{eq:prep-umbral}
\end{equation}

que combina dos criterios. La tensión que resuelve es la siguiente: una
lectura físicamente imposible hay que retirarla, pero un criterio
demasiado estricto retira consumo real, y mirando el valor aislado no hay
forma de distinguirlos.

El primer criterio mide cuánto se dispersa la mitad central de las
lecturas y pone el corte cinco veces esa dispersión por encima de ella.
Sigue la regla de Tukey, la que traza los bigotes de un diagrama de caja,
pero con un multiplicador de 5 en vez del 1,5 clásico, deliberadamente
permisivo para que solo retire lo inverosímil. Funciona bien mientras la
serie tenga recorrido.

Donde falla es en las cargas muy planas. Si la mitad central de las
lecturas cabe en un margen estrecho, el corte queda cerca de esa base y
cualquier pico operativo legítimo lo supera. CESMAG es ese caso: bajo M1 su
mitad central mide poco más de un kilovatio y el corte cae en
\uni{10.6}{kW}, un valor que la institución rebasa con frecuencia.

El segundo criterio existe para eso. Fija un valor por debajo del cual el
corte no puede bajar, y lo sitúa un \pct{20} por encima del percentil 99,5,
es decir, por encima de lo que la propia serie alcanza operando con
normalidad. Como el umbral es el mayor de los dos, se aplica siempre el
\emph{más permisivo}, y el segundo solo entra cuando el primero se ha vuelto
demasiado estricto.

La Figura~\ref{fig:prep-anatomia} construye los dos candidatos sobre las
diez series de demanda, de modo que la multiplicación por cinco se cuenta
con el ojo en vez de leerse en la fórmula.

\begin{figure}[H]
 \centering
 \includegraphics[width=0.95\textwidth]{figuras/f3_04b_anatomia_umbral.png}
 \caption[La anatomía del umbral]{Cómo se construye cada uno de los dos
 candidatos, en las diez series de demanda y en las dos fronteras. La caja
 oscura es la mitad central de las lecturas; sobre el tercer cuartil se
 apilan cinco copias exactas de esa misma caja, y donde aterriza la quinta
 está el primer criterio. A su derecha, la línea de puntos sube hasta el
 percentil 99,5 y el bloque ámbar es su prolongación del \pct{20}, cuyo
 techo es el segundo criterio. El que fija el umbral va con trazo lleno y
 el descartado, discontinuo.}
 \label{fig:prep-anatomia}
\end{figure}

\notafig{Cada panel lleva su propia escala vertical, porque los rangos van
de \uni{0.3}{kW} a \uni{24.4}{kW}, de modo que lo que se compara entre
paneles no es una altura sino una relación, esto es, si la cola con su
prolongación queda por encima o por debajo de la quinta copia. CESMAG bajo
M3 es el caso extremo: su pila apenas despega del suelo del panel y el
umbral queda cuatro veces más arriba, porque su cola cae a 23,2 rangos del
cuartil. El Hospital, con el rango más estrecho de todo el estudio,
\uni{0.318}{kW}, no necesita el segundo criterio en ninguna de las dos
fronteras. Los valores de cada panel están en la
Tabla~\ref{tab:prep-umbral}. Elaboración propia desde el caché de estados
intermedios.}

Lo que está en juego se mide en CESMAG. Sin ese segundo criterio, su
demanda perdería \num{451} horas bajo M1, que concentran el \pct{20.7} de
su energía, y \num{927} bajo M3, que concentran el \pct{52.5}. Con él tiene
cero atípicos bajo M1 y dos bajo M3.

El segundo criterio manda en cuatro de las 20 series, todas de demanda:
Mariana y CESMAG en las dos fronteras. En las 16 restantes basta el
primero, y entre ellas están las diez de generación, donde además no se
retira una sola hora. La Tabla~\ref{tab:prep-umbral} da la anatomía
completa de las series de demanda.

\begin{table}[H]
 \centering
 \caption{Los dos candidatos en cada serie de demanda, con el que se
 aplica en negrita. La columna de la cola dice a cuántos rangos
 intercuartílicos del tercer cuartil cae el percentil 99,5, y es la que
 explica el reparto: cuanto más lejos queda la cola del cuerpo, más se
 queda corto el primer criterio.}
 \label{tab:prep-umbral}
 \footnotesize
 \setlength{\tabcolsep}{5pt}
 \begin{tabular}{@{}l r r r r r@{}}
 \toprule
 \textbf{Institución} & \textbf{Rango} & \textbf{Cola} &
 \textbf{Primer} & \textbf{Segundo} & \textbf{Horas} \\
 & \textbf{(kW)} & \textbf{(rangos)} & \textbf{criterio (kW)} &
 \textbf{criterio (kW)} & \textbf{retiradas} \\
 \midrule
 \multicolumn{6}{@{}l}{\itshape Frontera M1} \\
 Udenar & \num{4,91} & \num{3,4} & \textbf{\num{33,4}} & \num{31,0} & \num{0} \\
 Mariana & \num{4,50} & \num{4,1} & \num{33,5} & \textbf{\num{35,1}} & \num{4} \\
 UCC & \num{24,35} & \num{1,3} & \textbf{\num{154,8}} & \num{78,5} & \num{0} \\
 HUDN & \num{1,40} & \num{1,2} & \textbf{\num{16,8}} & \num{13,7} & \num{0} \\
 CESMAG & \num{1,27} & \num{8,2} & \num{10,6} & \textbf{\num{17,7}} & \num{0} \\
 \addlinespace
 \multicolumn{6}{@{}l}{\itshape Frontera M3} \\
 Udenar & \num{1,54} & \num{2,2} & \textbf{\num{9,5}} & \num{6,3} & \num{0} \\
 Mariana & \num{0,97} & \num{5,0} & \num{7,6} & \textbf{\num{9,1}} & \num{5} \\
 UCC & \num{3,62} & \num{1,2} & \textbf{\num{23,0}} & \num{11,1} & \num{0} \\
 HUDN & \num{0,32} & \num{1,5} & \textbf{\num{3,7}} & \num{3,0} & \num{0} \\
 CESMAG & \num{0,56} & \num{23,2} & \num{4,5} & \textbf{\num{17,6}} & \num{2} \\
 \bottomrule
 \end{tabular}
\end{table}

Que la mitad central sea estrecha no lo decide por sí solo, y la tabla lo
enseña en dos filas contiguas: bajo M1 el HUDN y CESMAG tienen rangos casi
iguales y el umbral de cada uno lo fija un criterio distinto. Lo que los
separa es la columna de la cola, \num{1,2} rangos frente a \num{8,2}. Eso
es lo que obliga al segundo criterio, una base estable con picos altos y
poco frecuentes, y no la estrechez del cuerpo.

El criterio retira \num{11} horas en las 20 series. Todas son de demanda y
solo de dos instituciones: Mariana, con cuatro bajo M1 y cinco bajo M3, y
CESMAG, con dos bajo M3.

Queda una objeción por responder, y es que el multiplicador de 5 no
procede de ninguna referencia sino de la inspección de estas series. La
respuesta no es defender el valor sino medir cuánto depende de él el
resultado. La Tabla~\ref{tab:prep-sensibilidad} barre el multiplicador
entre el clásico y el doble del elegido, con los dos criterios y con el
primero solo.

\begin{table}[H]
 \centering
 \caption{Energía que el criterio retira, en porcentaje de la demanda de
 las 10 series, al variar el multiplicador del primer criterio.}
 \label{tab:prep-sensibilidad}
 \small
 \begin{tabular}{@{}c r r@{}}
 \toprule
 \textbf{Multiplicador} & \textbf{Con los dos} & \textbf{Solo con el} \\
 & \textbf{criterios} & \textbf{primero} \\
 \midrule
 1,5 & \pct{0.083} & \pct{13.369} \\
 3 & \pct{0.083} & \pct{6.563} \\
 \textbf{5} & \textbf{\pct{0.062}} & \pct{3.650} \\
 7 & \pct{0.021} & \pct{2.356} \\
 \bottomrule
 \end{tabular}
\end{table}

La columna de la izquierda contesta la objeción: entre el multiplicador
clásico y el elegido hay cuatro horas de diferencia sobre las \num{122880}
horas de serie del estudio, y la energía afectada no llega en ningún caso
al \pct{0.1} de la demanda. El valor del multiplicador no gobierna nada.

La columna de la derecha explica por qué. Sin el segundo criterio, el
multiplicador sí gobierna, y de forma brutal: el clásico retiraría el
\pct{13.4} de la demanda y ni siquiera un 7 bajaría del \pct{2.4}. El
segundo criterio no cubre un fallo del primero, sino que hace que el
primero deje de importar, y es esa insensibilidad, y no la elección del
valor, lo que sostiene el criterio.

La Figura~\ref{fig:prep-umbral} muestra a la izquierda qué le pasa a una de
esas horas y a la derecha por qué el umbral está donde está en cada serie.

\begin{figure}[H]
 \centering
 \includegraphics[width=0.95\textwidth]{figuras/f3_05_umbral_atipicos_m1.png}
 \caption[El umbral de atípicos y lo que retira]{A la izquierda, la hora
 que el umbral retira con mayor holgura entre las que quedan aisladas, la
 de la Universidad Mariana el 1 de octubre de 2025, con el valor que entra
 y el que queda tras la cascada. A la derecha, el criterio en las diez
 series, con el eje en múltiplos del percentil 99,5 de cada una: la barra
 es la mitad central de las lecturas, la línea llega hasta el máximo
 observado, el círculo marca el corte del primer criterio y el rombo el
 umbral que se aplica.}
 \label{fig:prep-umbral}
\end{figure}

\notafig{La hora del caso entra con \uni{39.9}{kW}, excede en un
\pct{13.7} el umbral de \uni{35.1}{kW} y sale con \uni{28.2}{kW},
interpolada entre las lecturas de las 8 y las 10 horas. En el panel
derecho, el segundo criterio es la vertical en 1,2, común a las diez filas,
de modo que manda justamente donde el círculo cae a su izquierda. La
distancia entre la masa de la serie y su umbral es lo que acredita el
criterio, y solo en cuatro de las 20 series hace falta el segundo para
conservarla. Elaboración propia desde el caché de estados intermedios.}

\subsection{Tratamiento de los huecos}
\label{sub:prep-huecos}

Lo que queda como faltante, sea porque el equipo no registró o porque el
umbral acaba de retirarlo, pasa por una cascada de cuatro tratamientos, de
los que solo tres llegan a actuar. La cascada no clasifica los huecos por
su longitud: se aplica hora a hora dentro de cada uno, de modo que un
mismo hueco puede recibir varios.

Las tres primeras horas de cualquier hueco se interpolan linealmente en el
tiempo, sobre la recta que une la última hora observada antes del hueco
con la primera de después. Si el hueco dura tres horas o menos, ahí
termina. Si es más largo, el tramo interpolado avanza solo esas tres horas
y el resto pasa al tratamiento siguiente.

Ese resto se rellena repitiendo un valor, primero hacia adelante hasta 24
horas y después hacia atrás otras 24. Hacia adelante, el valor que se
repite no es el que cabría esperar: se propaga la tercera hora
interpolada, que es un valor fabricado, y no la última lectura del equipo.
Hacia atrás sí se propaga una lectura real, la primera que reabre la
serie. En los dos casos se repite un único número, de modo que el tramo
queda plano en vez de recuperar el ciclo diario.

Sumados los tres tramos, tres horas más 24 más otras 24, la cascada cierra
cualquier hueco de hasta 51 horas. El cuarto tratamiento imputaría con
cero lo que sobreviviera, y por eso no llega a actuar en este horizonte:
el hueco más largo mide 49 horas, en la demanda de la UCC bajo M3.

La limpieza trata en total \num{808} horas repartidas entre las 20 series,
que son las \num{11} que el umbral retira más las \num{797} que ya
faltaban en el dato. De ellas, \num{233} quedan interpoladas y \num{575},
el \pct{71.2}, rellenas repitiendo un valor.

Que ese relleno sea plano tiene una consecuencia que hay que acotar, y es
que introduce mesetas donde el capítulo va a leer después la coincidencia
entre el pico solar y el valle de demanda. La cota es doble. La primera es
que \emph{la generación no se rellena ni una sola hora} en ninguna de las
dos fronteras, de modo que el pico solar nunca se fabrica: lo que la
limpieza toca es siempre demanda. La segunda es el tamaño de lo tocado.
Sobre las \num{30720} horas de serie de cada frontera, la limpieza rellena
\num{234} bajo M1 y \num{574} bajo M3, y de ellas caen en la franja diurna
\num{121} y \num{300}, es decir, el \pct{0.39} y el \pct{0.98} de las
horas de serie.

La Figura~\ref{fig:prep-huecos} recorre los tres tratamientos sobre huecos
reales de longitud creciente.

\begin{figure}[H]
 \centering
 \includegraphics[width=0.95\textwidth]{figuras/f3_06_escalera_huecos.png}
 \caption[La escalera de los huecos]{Tres huecos reales de longitud
 creciente, con el eje horizontal en horas desde el inicio de cada uno, que
 es la magnitud de la que depende el tratamiento, y cada panel con su
 frontera de medición, porque los tres no salen de la misma. Abajo, cuántas
 horas trató la limpieza en cada institución y en cada frontera, y por qué
 mecanismo.}
 \label{fig:prep-huecos}
\end{figure}

\notafig{El hueco de tres horas lo cierra entero la interpolación. El de 13
sale con sus tres primeras horas interpoladas y las diez restantes
arrastradas, y la meseta no se queda en la última lectura observada,
\uni{9.9}{kW}, sino en la tercera hora interpolada, \uni{10.1}{kW}, es
decir, en un valor que ninguna hora observada tuvo. El de 49, el más largo
de todo el estudio, añade el arrastre hacia atrás desde la lectura que
reabre la serie. Los ceros de M1 en Udenar y en la Universidad Cooperativa
de Colombia no significan que no tuvieran cortes: su medidor es neto y el
recorte de la reconstrucción cerró esos huecos antes de que la limpieza los
viera, y bajo M3 esas mismas dos instituciones acumulan \num{97} y
\num{126} horas. El arrastre es además el \pct{71.2} de todo lo que la
limpieza trata, de modo que la mayor parte de lo rellenado es una meseta y
no una rampa. Elaboración propia desde el caché de estados intermedios.}

\begin{trampabox}
En el proyecto conviven tres criterios de relleno distintos:

\begin{enumerate}[leftmargin=*, itemsep=0.2em, topsep=0.3em]
 \item Las 10 series de medida interpolan las tres primeras horas de cada
 hueco, arrastran el resto hasta 24 horas por cada lado y solo entonces
 recurrirían al cero.
 \item La serie horaria de precios de bolsa, que construye el
 Capítulo~\ref{sec:precios}, interpola las seis primeras horas de cada
 hueco y completa el resto con la mediana de la propia serie, sin marcado
 de atípicos y sin arrastre.
 \item El caso de estudio prospectivo interpola también las seis primeras
 y completa el resto con cero.
\end{enumerate}

Las diferencias son tres y solo una está justificada. Que el residuo de
los precios sea la mediana y no cero se sostiene, porque un consumo de
cero kilovatios es una lectura legítima pero un precio de cero no describe
ninguna hora del mercado mayorista. El límite de interpolación, seis horas
en vez de tres, y la ausencia de arrastre en los dos últimos criterios no
tienen justificación registrada: son heterogeneidad heredada de la
construcción incremental del repositorio, y no elecciones fundamentadas.
\end{trampabox}

\subsection{Ensamblaje y verificación de las matrices}
\label{sub:prep-matrices}

Las 10 series limpias se apilan en dos matrices de dimensión
$5 \times \num{6144}$ y se someten a una verificación de contrato: si
alguna celda de $D$ o de $G$ resulta negativa, la ejecución se detiene con
un error explícito en lugar de continuar con datos inconsistentes. La
no negatividad no es aquí una aspiración sino una propiedad garantizada
por construcción y comprobada al final. La
Figura~\ref{fig:prep-matrices} despliega las dos matrices como mapas de
hora contra día, que es la forma en que sus estructuras se vuelven
visibles de una vez.

\begin{figure}[H]
 \centering
 \includegraphics[width=0.95\textwidth]{figuras/f3_07_matrices_m1.png}
 \caption[Las matrices que consume el modelo]{Las dos matrices
 resultantes, dibujadas como mapas de hora del día contra día del
 horizonte, en la frontera M1. La generación traza una banda diurna
 nítida que sigue el arco solar; la demanda muestra la semana laboral
 en franjas verticales regulares.}
 \label{fig:prep-matrices}
\end{figure}

\begin{lecturabox}
Esta figura funciona como comprobación final del capítulo. Si hubiera un
error de zona horaria, la banda solar aparecería desplazada respecto del
mediodía; si la alineación temporal entre medidores e inversores fallara,
la banda se difuminaría; si la agregación horaria estuviera mal, las
franjas de la semana laboral no se distinguirían. Que las tres estructuras
aparezcan donde deben es lo que autoriza a seguir adelante.
\end{lecturabox}

\subsection{La zona horaria del índice}
\label{sub:prep-zona}

La última etapa no altera ningún valor: le pone zona horaria al índice.
Aun así encierra una suposición de la que dependen dos argumentos
anteriores.

Las marcas de tiempo llegan sin zona, de modo que hay que decidir cuál es
la suya. El pipeline las declara hora local de Colombia, es decir, las
interpreta como ya escritas en ese huso en vez de convertirlas desde otro.
La suposición es que cada equipo escribe la hora del sitio donde está
instalado, y se comprueba sin salir del dato: el máximo de generación de
las cinco instituciones cae entre las 11 y las 12, que es el mediodía solar
de Pasto. Si las marcas vinieran en tiempo universal, con Colombia cinco
horas por detrás, ese máximo aparecería entre las 16 y las 17.

Los dos argumentos que dependen de esto son el de la
subsección~\ref{sub:prep-negativa}, que sitúa las lecturas negativas en las
horas de sol, y el de la subsección~\ref{sub:prep-matrices}, que lee la
banda solar de la matriz de generación en torno al mediodía. Los dos se
caerían si la serie estuviera desplazada cinco horas.

Colombia no aplica horario de verano, de modo que en este conjunto no hay
horas que no existan ni horas que ocurran dos veces. Las dos salvaguardas
que el pipeline lleva para esos casos, una que adelanta la hora inexistente
y otra que resuelve la ambigua, nunca llegan a actuar.

\subsection{Cinco ritmos distintos}
\label{sub:prep-perfiles}

Con las matrices ya construidas se puede mirar lo que caracteriza a esta
comunidad y hace posible el intercambio: las cinco instituciones comparten
municipio, clima y recurso solar, pero no comparten rutina. La
Figura~\ref{fig:prep-perfiles} pone los cinco perfiles uno junto a otro,
con el de la comunidad completa como sexto panel.

\begin{figure}[H]
 \centering
 \includegraphics[width=0.95\textwidth]{figuras/f3_08_perfiles_instituciones_m1.png}
 \caption[Perfil medio de cada institución]{Perfil medio de demanda y
 generación por institución en la frontera M1; el porcentaje del
 título es la razón entre generación y demanda de cada una. La curva
 verde sombreada es común en forma a las cinco (el sol es el mismo),
 mientras que las curvas de demanda difieren en nivel y en forma.}
 \label{fig:prep-perfiles}
\end{figure}

\notafig{El HUDN, que es un hospital, mantiene una demanda casi plana
las 24 horas, como corresponde a una instalación que no cierra.
Udenar dibuja el doble pico característico de un campus universitario,
con el valle del almuerzo entre ambos. La UCC opera al nivel más alto del
conjunto en la frontera M1 y concentra su consumo en la jornada diurna.
Elaboración propia desde el caché de estados intermedios.}

Esa disparidad de formas es exactamente la condición que hace que unas
instituciones tengan excedente cuando otras tienen déficit, y sin ella no
habría nada que intercambiar.

\subsection{El ritmo semanal y el anual}
\label{sub:prep-ritmos}

Hay una segunda fuente de disparidad, esta vez temporal y común a las
cinco: el calendario. La Figura~\ref{fig:prep-ritmos} cruza el perfil
diario con el corte semanal y con el recorrido mensual.

\begin{figure}[H]
 \centering
 \includegraphics[width=0.95\textwidth]{figuras/f3_09_ritmos_m1.png}
 \caption[Ritmo semanal y recorrido anual]{A la izquierda, día hábil
 frente a fin de semana: la demanda cae de forma apreciable mientras la
 generación permanece igual. A la derecha, el recorrido mensual de las
 dos magnitudes a lo largo del horizonte.}
 \label{fig:prep-ritmos}
\end{figure}

\notafig{El corte semanal esconde una asimetría que reaparece en los
resultados: el fin de semana la demanda cae pero la generación no, porque
el sol no distingue entre sábado y martes. La consecuencia es que el
excedente disponible para intercambio es mayor precisamente cuando hay
menos gente en los edificios, es decir, cuando la comunidad tiene menos
capacidad de absorberlo internamente. Elaboración propia desde el caché
de estados intermedios.}

\subsection{Las decisiones que sesgan en contra}
\label{sub:prep-sesgo}

En las seis etapas se decidió algo y el resultado sale favorable, de modo
que la objeción es inmediata: cuánto del resultado son las decisiones. La
respuesta honesta no es que las decisiones sean neutrales, sino que lo que
importa es hacia dónde apuntan.

Seis apuntan en la misma dirección, la de recortar el resultado. Cuatro son
del pipeline y dos del modelo:

\begin{enumerate}[leftmargin=*, itemsep=0.35em]
 \item \textbf{El recorte de la reconstrucción} elimina
 \uni{1013.3}{kWh} de la demanda de Udenar en horas solares, según la
 Tabla~\ref{tab:prep-reconstruccion}. Menos demanda en horas de sol
 significa menos déficit que cubrir y, por tanto, un mercado más
 pequeño.

 \item \textbf{Un solo medidor y un solo inversor por institución.} El
 pipeline no suma todos los equipos disponibles sino que designa uno de
 cada clase. En Udenar, que aporta tres inversores y expone uno, la
 generación que ve el mercado queda por debajo de la instalada; en las
 otras cuatro el equipo designado es el único, de modo que coinciden.
 En el agregado comunitario, por tanto, se expone menos de lo que hay.

 \item \textbf{La ausencia de gestión de la demanda.} El modelo del
 Capítulo~\ref{sec:modelo} corre con la demanda tal como se midió, sin
 programa de desplazamiento de carga. Toda ganancia que un despliegue
 real obtendría moviendo consumo hacia las horas de sol queda fuera del
 resultado.

 \item \textbf{El descarte conservador de las horas numéricamente
 difíciles.} Cuando el método numérico que resuelve el equilibrio
 de una hora no converge, esa hora se descarta sin imputarle beneficio
 al mercado entre pares, en vez de estimarlo.

 \item \textbf{La generación que falta se pone a cero antes de la
 limpieza}, de modo que sus huecos no se interpolan ni se arrastran como
 los de la demanda. En la noche eso es correcto, porque el equipo no
 produce, pero \num{285} horas de sol del horizonte quedan con generación
 nula cuando lo que hubo fue una avería de registro. Menos generación es
 menos excedente que ofrecer.

 \item \textbf{La demanda sin lectura de medidor también se pone a cero},
 en la reconstrucción y no en la limpieza, según la
 subsección~\ref{sub:prep-reconstruccion}. Son \num{330} horas bajo M1 que
 entran al modelo como consumo nulo, y menos demanda es menos déficit que
 cubrir.
\end{enumerate}

La séptima apunta al revés. Al estimar la hora incompleta con la media de
sus muestras presentes se le atribuye al tramo no observado el
comportamiento del observado, lo que da entre \pct{0.26} y \pct{0.55} más
energía que truncarlo. Afecta a \num{835} horas del horizonte, es decir, al
\pct{2.8} de las horas con dato, en las que se observa de media el
\pct{85.9} de la hora.

No se elige por su signo: truncar queda desmentido por el contador del
propio equipo y va en contra del criterio del regulador, según la
subsección~\ref{sub:prep-incompleta}. El supuesto se comprobó además
aplicando la ventana observada de cada hora incompleta a horas completas
comparables de la misma institución, la misma hora del día y el mismo mes:
el sesgo resulta indistinguible de cero en las cinco, y suma el
\pct{0.0007} de la demanda del horizonte.

Seis de las siete recortan el resultado y la restante lo amplía en una
proporción tres órdenes de magnitud menor que el margen que separa a los
mecanismos. La
consecuencia metodológica es que las cifras del Capítulo~\ref{sec:res}
deben leerse como una \textbf{cota inferior} de lo que este mercado podría
rendir, y no como una estimación centrada. Es una posición deliberadamente incómoda, porque renuncia a parte del
margen, y es la que permite sostener el resultado ante una lectura
hostil.

\subsection{Declaración de honestidad metodológica}
\label{sub:prep-honestidad}

Cuatro elecciones de este capítulo se declaran sin atenuarlas. Las tres
primeras carecen de respaldo formal; la cuarta está adoptada pero todavía
no aplicada.

\begin{enumerate}[leftmargin=*, itemsep=0.35em]
 \item El umbral de valores atípicos
 $\max(Q_{75} + 5\,\mathrm{IQR},\; 1{,}2 \cdot P_{99{,}5})$ se adoptó
 por inspección de las series del proyecto, y no hay referencia de
 literatura que respalde el multiplicador cinco. Lo que sí hay es la
 medida de cuánto depende el resultado de él, según la
 Tabla~\ref{tab:prep-sensibilidad}: entre el multiplicador clásico y el
 elegido la energía retirada varía dos centésimas de punto porcentual, de
 modo que la falta de respaldo bibliográfico no compromete la cifra. Lo
 que sostiene el criterio no es el valor sino esa insensibilidad.

 \item El relleno hacia adelante con límite de 24 horas es una
 heurística de «día operativo» (supone que el patrón del día anterior es
 representativo) cuyo impacto no se evaluó de forma aislada.

 \item No se construyó una matriz de sensibilidad del pipeline completo,
 es decir, variar el umbral y los límites de interpolación y relleno y
 medir el efecto sobre los resultados. Queda como trabajo futuro.
 \item Se adopta como criterio de validez que una hora traiga al menos el
 \pct{75} de sus muestras, y la corrida que este documento reporta es
 anterior a esa adopción, de modo que no lo aplica: las \num{147} horas
 que no lo cumplen entraron estimadas como las demás. El efecto está
 medido y vale el \pct{0.026} de la demanda comunitaria, según la
 subsección~\ref{sub:prep-incompleta}. Incorporarlo exige rehacer la corrida
 canónica completa, de modo que queda pendiente para la próxima, y hasta
 entonces la cifra de arriba es la cota de lo que cambiaría.
\end{enumerate}

Como cota global, el impacto neto de estas elecciones sobre el volumen de
energía transada se estima en $\pm$\,\pct{0.3}. El estatuto de esa cifra
necesita una precisión, porque es fácil leerla como un resultado y no lo
es: no procede de la matriz de sensibilidad que el
punto~3 declara no construida, sino de un juicio apoyado en dos evidencias indirectas: la cobertura de
adquisición por equipo, que supera el \pct{97} en los inversores
principales, y la magnitud del recorte de la reconstrucción, que
representa el \pct{0.32} de la demanda comunitaria. Es una estimación
declarada y no un valor calculado, y así debe leerse.

Con todo lo anterior, el dato queda listo para el modelo: dos matrices de
\num{6144} horas, no negativas por construcción, con la energía que cada
institución consume y la que produce separadas de forma explícita, y con
el costo de cada intervención cuantificado. Queda, sin embargo, una decisión
que este capítulo ha mencionado sin desarrollar, la de qué medidor se lee
en cada institución, y de ella depende buena parte de lo que sigue: es el
objeto del capítulo siguiente.

% =====================================================================
% La frontera de medición: M1 y M3
% Parte I. El dato
%
% ENCARGO: dos coberturas, dos preguntas. La bifurcacion es en si misma
% un resultado: invierte los papeles de los agentes.
%
% Reglas del documento:
% - Una idea = una subseccion = una figura.
% - Toda figura declara su procedencia con \fuente{...}.
% - Toda cifra sale del canon de agosto; ver GUIA.md.
% - M1 y M3 se reportan siempre en paralelo.
% - Ningun superlativo sin releer el CSV de la frontera correcta.
% =====================================================================
\section{La frontera de medición: M1 y M3}
\label{sec:frontera}
\justifying

El capítulo anterior terminó con dos matrices y una advertencia: la
bifurcación entre las coberturas M1 y M3 ocurre en la primera etapa del
pipeline y en ninguna otra. Este capítulo desarrolla esa bifurcación, y lo
hace en capítulo aparte por una razón que conviene anticipar: no se trata
de dos versiones del mismo cálculo, ni de un análisis de sensibilidad
sobre un parámetro. Se trata de dos preguntas distintas, y las dos tienen
respuesta. Elegir una sola y presentarla como \emph{el} resultado sería la
decisión más discutible que este trabajo podría tomar.

\subsection{Dos preguntas, no dos versiones}
\label{sub:frontera-preguntas}

Cada institución de la comunidad tiene más de un medidor, y el
inventario de instalación del Capítulo~\ref{sec:dato} estableció qué
circuito registra cada uno: el Medidor 1, el circuito principal o de
inyección; el Medidor 3, un circuito secundario. Ninguno totaliza la
institución completa ni aísla el ramal que alimenta el sistema
fotovoltaico. Leer uno u otro no es una preferencia técnica:
define contra qué se compara la generación, y por tanto qué pregunta se
está respondiendo. La Figura~\ref{fig:frontera-esquema} ubica las dos
fronteras sobre el circuito de una institución.

\begin{figure}[H]
 \centering
 \begin{tikzpicture}[
 font=\scriptsize,
 blk/.style={draw=black!55, fill=black!3, rounded corners=1.5pt,
 minimum width=1.75cm, minimum height=0.6cm,
 align=center},
 med/.style={circle, draw=black!70, fill=white, minimum size=0.62cm,
 inner sep=0pt, font=\scriptsize\bfseries},
 inv/.style={draw=successcolor, fill=successcolor!14, very thick,
 rounded corners=1.5pt, minimum width=1.35cm,
 minimum height=0.55cm, align=center},
 lin/.style={-{Stealth[length=3.5pt]}, draw=black!60},
 et/.style={font=\scriptsize\itshape, align=left},
 ]
 \node[blk] (red) at (0,0) {Red};
 \node[med, draw=sectioncolor, fill=sectioncolor!14, very thick]
 (m1) at (2.0,0) {M1};
 \node[blk, minimum width=2.1cm] (campus) at (4.4,0)
 {Circuito principal\\o de inyección};
 \node[inv] (inv) at (7.1,0) {Inversor};

 \node[med, draw=warningcolor, fill=warningcolor!16, very thick]
 (m3) at (4.4,-1.35) {M3};
 \node[blk, minimum width=2.1cm] (circ) at (4.4,-2.55)
 {Circuito\\secundario};

 \draw[lin] (red) -- (m1);
 \draw[lin] (m1) -- (campus);
 \draw[lin] (inv) -- (campus);
 \draw[lin] (campus) -- (m3);
 \draw[lin] (m3) -- (circ);

 \node[et, text=sectioncolor, right=0.5cm of inv, anchor=west,
 text width=3.9cm]
 {\textbf{M1} mide el circuito principal\\
 o de inyección: la generación\\
 equivale al 19,1 \% de su consumo};
 \node[et, text=warningcolor, right=0.5cm of circ, anchor=west,
 text width=3.9cm]
 {\textbf{M3} mide un circuito\\
 secundario: la generación\\
 equivale al 91,2 \% de su consumo};
 \end{tikzpicture}
 \caption[Las dos fronteras de medición]{Posición de las dos fronteras
 de medición en el circuito de una institución. El mismo inversor
 entrega la misma energía en los dos casos; lo que cambia es la carga
 contra la que se la compara. El esquema es una idealización: la posición
 del punto de inyección respecto de cada medidor varía entre
 instituciones, como recoge el inventario del Capítulo~\ref{sec:dato}.
 \fuente{elaboración propia a partir de \code{DEMAND\_METER\_CONFIG} y
 \code{PAPER\_METER\_DEMAND\_CONFIG} en \file{data/preprocessing.py}.}}
 \label{fig:frontera-esquema}
\end{figure}

La cobertura \textbf{M1} pregunta: \emph{¿qué le aporta este mercado a
una carga grande, frente a la cual la generación es escasa?} La cobertura
\textbf{M3} pregunta: \emph{¿qué le aporta a una carga del orden de la
propia generación?} La primera describe un mercado con poco excedente que
repartir; la segunda, uno donde el excedente abunda. Ninguna es más
legítima que la otra, porque la relación de cada circuito con el consumo
institucional completo está por establecer.

\subsection{Qué mide cada frontera}
\label{sub:frontera-cobertura}

La diferencia entre las dos coberturas se resume en un hecho que conviene
subrayar porque desactiva de entrada una posible objeción: \textbf{la
generación es idéntica en las dos}. El pipeline designa el mismo inversor
por institución en un caso y en el otro, de modo que la matriz $G$ es la
misma matriz. Lo único que cambia es el denominador. La
Figura~\ref{fig:frontera-cobertura} confronta las dos energías en el
agregado comunitario y en cada institución.

\begin{figure}[H]
 \centering
 \includegraphics[width=0.95\textwidth]{figuras/f4_02_cobertura.png}
 \caption[Cobertura de generación en cada frontera]{Energía medida en
 el horizonte bajo cada frontera. A la izquierda, la comunidad
 completa: la barra de generación es la misma en los dos casos y solo
 cambia la de demanda. A la derecha, la razón entre generación y
 demanda por institución. Las cifras se leen del caché del
 preprocesamiento, no se fijan en el código.
 \fuente{elaboración propia desde el caché de estados intermedios del
 preprocesamiento.}}
 \label{fig:frontera-cobertura}
\end{figure}

\notafig{La demanda comunitaria medida es de \uni{318046}{kWh} bajo M1
y de \uni{66734}{kWh} bajo M3, mientras que la generación es de
\uni{60860}{kWh} en ambos casos, es decir, exactamente la misma cifra. De
ahí las dos coberturas, \pct{19.1} y \pct{91.2}, que no reflejan dos
instalaciones distintas sino dos maneras de delimitar la carga.
Elaboración propia desde el caché de estados intermedios.}

\begin{lecturabox}
Conviene resistir la lectura de que M3 «tiene más generación». No la
tiene: tiene la misma. Lo que M3 hace es comparar esa generación contra
una carga cinco veces menor, porque solo mira un circuito secundario de
cada institución. Una generación que equivale al \pct{19} del consumo del
circuito principal puede equivaler al \pct{91} del consumo de un circuito
menor sin que se haya instalado un solo panel adicional.
\end{lecturabox}

\subsection{Los mismos días, dos lecturas}
\label{sub:frontera-perfiles}

El mismo hecho visto en el eje del tiempo resulta más elocuente que en
totales. La Figura~\ref{fig:frontera-perfiles} reúne el perfil diario de
las dos coberturas en un mismo par de paneles.

\begin{figure}[H]
 \centering
 \includegraphics[width=0.95\textwidth]{figuras/f4_03_perfiles_m1_m3.png}
 \caption[Perfil comunitario bajo cada frontera]{Perfil medio por hora
 del día de la comunidad en cada cobertura. La curva verde es la misma
 en los dos paneles; la gris cae casi un factor cinco. En M3 las dos
 curvas llegan a cruzarse en las horas centrales, lo que en M1 no
 ocurre nunca.
 \fuente{elaboración propia desde el caché de estados intermedios del
 preprocesamiento.}}
 \label{fig:frontera-perfiles}
\end{figure}

\notafig{El cruce del panel derecho es la consecuencia operativa de
todo el capítulo. Bajo M1 la generación no alcanza a la demanda
comunitaria en ninguna hora del día, de modo que la comunidad es siempre
compradora neta frente a la red y el intercambio interno se limita a
redistribuir un excedente escaso. Bajo M3 hay horas en que la comunidad
genera más de lo que su circuito fotovoltaico consume, y el excedente
disponible para intercambio es sustancialmente mayor. Elaboración propia
desde el caché de estados intermedios.}

\subsection{La inversión de los papeles}
\label{sub:frontera-inversion}

Hasta aquí el efecto es de escala, y podría parecer que M3 es simplemente
«M1 con más recurso relativo». No lo es. El cambio de denominador altera
qué institución tiene excedente y cuál tiene déficit en cada hora, y el
resultado es que los papeles se invierten. La
Figura~\ref{fig:frontera-inversion} cruza las dos coberturas conservando
el orden de las instituciones, de modo que el vuelco se lee de un panel a
otro.

\begin{figure}[H]
 \centering
 \includegraphics[width=0.95\textwidth]{figuras/f4_04_inversion_papeles.png}
 \caption[Reparto de la energía transada en cada frontera]{Participación
 de cada institución en la energía transada, hacia la derecha como
 vendedora y hacia la izquierda como compradora. Las cinco conservan el
 mismo orden vertical en los dos paneles, de modo que el cambio de
 frontera se lee como un vuelco del perfil.
 \fuente{elaboración propia desde
 \file{p2p\_breakdown\_flujos.csv} de las dos carpetas canónicas.}}
 \label{fig:frontera-inversion}
\end{figure}

\notafig{Bajo M1, Udenar aporta el \pct{76.3} de la energía vendida y
la UCC absorbe el \pct{36.3} de la comprada. Bajo M3, Udenar baja al
\pct{30.5} de lo vendido, el HUDN pasa al primer lugar con el \pct{33.9},
y Cesmag, que en M1 vendía el \pct{8.0}, absorbe el \pct{70.9} de todo lo
comprado. No es un reordenamiento menor: es un cambio de papel.
Elaboración propia desde \file{p2p\_breakdown\_flujos.csv} de las dos
carpetas canónicas.}

\begin{trampabox}
Esta figura es el origen de la regla de redacción más estricta del
documento. Ninguna frase con superlativo («el mayor vendedor», «solo dos
instituciones», «ninguna») puede escribirse sin decir de qué cobertura
habla, porque la institución que domina la venta en una frontera apenas
participa en la otra. Toda afirmación de ese tipo en los capítulos
siguientes se comprobó releyendo el archivo de flujos de la frontera
correspondiente.
\end{trampabox}

\subsection{El embudo de horas de mercado}
\label{sub:frontera-embudo}

Queda una magnitud que ninguna de las figuras anteriores muestra y que
conviene fijar antes de llegar a los resultados: de las \num{6144} horas
del horizonte, en muy pocas llega a haber mercado. La
Figura~\ref{fig:frontera-embudo} presenta esa reducción en cuatro
escalones sucesivos.

\begin{figure}[H]
 \centering
 \includegraphics[width=0.95\textwidth]{figuras/f4_05_embudo_horas.png}
 \caption[Del horizonte a las horas con mercado]{Reducción sucesiva del
 número de horas: horizonte completo, horas con generación, horas en
 que alguna institución tiene excedente sobre su propia demanda, y
 horas en que efectivamente se liquidó un intercambio. El último
 escalón procede de los archivos de flujos y no de una estimación.
 \fuente{elaboración propia desde el caché del preprocesamiento y
 \file{p2p\_breakdown\_flujos.csv} de las dos carpetas canónicas.}}
 \label{fig:frontera-embudo}
\end{figure}

\notafig{El recorrido es el mismo en las dos coberturas hasta el
segundo escalón, porque hay generación en \num{3353} horas y esa cifra no
depende de la frontera. La separación empieza en el tercero: bajo M1
quedan \num{1047} horas con excedente y se transa en \num{1029}; bajo M3,
\num{2512} y \num{1752}. Elaboración propia desde el caché del
preprocesamiento y los archivos de flujos.}

\begin{lecturabox}
Que el mercado opere en el \pct{16.7} de las horas bajo M1 y en el
\pct{28.5} bajo M3 no es un defecto del mecanismo: es una consecuencia
del recurso. Un mercado de excedentes solares no puede funcionar de noche.
Lo que la figura acota es la ambición razonable del ejercicio: cualquier
ventaja que el Capítulo~\ref{sec:res} reporte se genera en esa fracción de
horas y se diluye sobre el total.
\end{lecturabox}

\subsection{Una salvedad que hay que declarar}
\label{sub:frontera-mariana}

La cobertura M3 no es perfectamente homogénea, y conviene decirlo aquí y
no en una nota al pie. Cuatro de las cinco instituciones tienen un
medidor de circuito secundario que el pipeline lee directamente. Mariana
no lo tiene. Para representarla en M3, el pipeline
reutiliza su medidor de circuito principal escalado por un factor de $0{,}3$, bajo
el supuesto de que el circuito relevante equivale aproximadamente a una
facultad.

Ese factor es una convención, no una medición, y se declara como tal.
Afecta a una de las cinco instituciones y solo en una de las dos
coberturas. Su efecto sobre el resultado agregado se examina en el
Capítulo~\ref{sec:rob}, donde el barrido del nivel de demanda cubre
holgadamente el rango en que ese factor podría equivocarse.

\subsection{La consecuencia: reportar dos veces}
\label{sub:frontera-consecuencia}

De todo lo anterior se sigue la decisión metodológica que gobierna el
resto del documento: \textbf{cada resultado se reporta bajo las dos
coberturas, siempre, y sin elegir una como principal.}

La tentación de elegir es real y conviene nombrarla. M1 es la frontera en
la que el mercado entre pares obtiene su mejor resultado relativo, y una
presentación que se quedara solo con M1 mostraría un trabajo más
concluyente. M3 es la frontera en la que la generación pesa más sobre
la carga medida, con una razón entre generación y demanda del
\pct{91.2}; y es también la frontera en la que, como se verá, el
mecanismo colectivo supera el mercado entre pares. Presentar solo una
de las dos sería, en el primer caso, seleccionar el resultado favorable, y
en el segundo, ocultar el caso en que el mecanismo propuesto funciona
mejor.

La convención adoptada evita las dos: las figuras de los capítulos de
resultados llevan siempre dos paneles, el izquierdo M1 y el derecho M3, y
cuando las dos coberturas dan lecturas opuestas, se dicen las dos y se
discute por qué difieren.

En síntesis, la bifurcación deja dos preguntas con respuesta y una regla
de reporte: dos paneles por figura, con una razón entre generación y
demanda del \pct{19.1} bajo M1 y del \pct{91.2} bajo M3. Con la energía
así delimitada queda pendiente la otra mitad del dato, los precios, que el
capítulo siguiente construye desde sus fuentes primarias.

% =====================================================================
% La otra mitad del dato: los precios
% Parte I — El dato
%
% ENCARGO: bolsa XM, techo de escasez, tarifa CEDENAR y su desglose,
% LCOE solar y la banda de precios que resulta.
%
% Reglas del documento:
% - Una idea = una subseccion = una figura.
% - Toda figura declara su procedencia con \fuente{...}.
% - Toda cifra sale del canon de agosto; ver GUIA.md.
% - M1 y M3 se reportan siempre en paralelo.
% - Ningun superlativo sin releer el CSV de la frontera correcta.
% =====================================================================
\section{La otra mitad del dato: los precios}
\label{sec:precios}
\justifying

Los tres capítulos anteriores construyeron una pareja de matrices de
energía. Con ellas todavía no se puede decir nada económico, porque un
kilovatio-hora no tiene un valor único: vale una cosa si la institución lo
compra a su comercializador, otra si lo inyecta a la red, y otra distinta
si se lo compra a la institución vecina. Este capítulo reconstruye esa
segunda mitad del dato desde las fuentes primarias, y termina delimitando
la banda dentro de la cual puede existir un intercambio que convenga a las
dos partes.

Conviene anticipar por qué el capítulo es más largo de lo que parecería
necesario. Todas las cifras del Capítulo~\ref{sec:res} son productos de
una energía por un precio. La energía se midió; el precio se construyó. Un
error de un diez por ciento en la construcción del precio es un error de
un diez por ciento en el resultado, y por eso cada componente del precio
se documenta con su fuente.

\subsection{El precio al que la red compra}
\label{sub:precios-bolsa}

Cuando una institución produce más de lo que consume y no encuentra
comprador dentro de la comunidad, el excedente va a la red y se liquida al
precio del mercado mayorista. Esa serie es, por tanto, el suelo económico
de todo el ejercicio: ningún vendedor aceptará dentro de la comunidad un
precio inferior al que obtendría fuera. La
Figura~\ref{fig:precios-bolsa} presenta la serie completa junto a su
patrón intradiario.

\begin{figure}[H]
 \centering
 \includegraphics[width=0.95\textwidth]{figuras/f5_01_bolsa.png}
 \caption[Serie horaria del precio de bolsa]{Precio del mercado
 mayorista sobre el horizonte del estudio. A la izquierda, la serie
 horaria en gris con la media diaria superpuesta; a la derecha, el
 patrón intradiario con su recorrido intercuartílico.
 \fuente{elaboración propia desde
 \file{data/precios\_bolsa\_xm\_api.csv}, recortado a las \num{6144}
 horas del horizonte.}}
 \label{fig:precios-bolsa}
\end{figure}

\notafig{La media del horizonte es de \uni{182.5}{COP/kWh} y la mediana
de \uni{114.3}{COP/kWh}. La distancia entre ambas no es un detalle
estadístico: indica una serie muy asimétrica, con la masa concentrada en
valores bajos y una cola larga de horas caras que llega a
\uni{2224.0}{COP/kWh}. Elaboración propia desde
\file{data/precios\_bolsa\_xm\_api.csv}.}

\begin{trampabox}
El archivo de caché de precios cubre hasta enero de 2026, mientras que el
horizonte del estudio termina el 16 de diciembre de 2025. Promediar el
archivo completo da \uni{193.6}{COP/kWh} en lugar de \uni{182.5}{}, una
diferencia del \pct{6.1} que se propagaría a toda comparación con la
tarifa. Todas las figuras y cifras de este capítulo usan la serie
recortada al horizonte; el recorte lo aplica el propio lector de datos,
para que no dependa de que cada script lo recuerde.
\end{trampabox}

\subsection{El precio al que la red vende}
\label{sub:precios-cu}

En el otro extremo está lo que la institución paga por la energía que
compra. En Colombia ese precio es el costo unitario, una suma de siete
componentes fijadas por resolución, y su composición importa aquí más de
lo que suele importar en un estudio energético: cada uno de los mecanismos
regulatorios del Capítulo~\ref{sec:esc} exime al usuario de componentes
distintas. Sin el desglose, esas diferencias no se pueden ni enunciar. La
Figura~\ref{fig:precios-cu} descompone el costo unitario mes a mes y
ordena sus componentes por peso.

\begin{figure}[H]
 \centering
 \includegraphics[width=0.95\textwidth]{figuras/f5_03_cu_desglose.png}
 \caption[Composición del costo unitario]{A la izquierda, el costo
 unitario mes a mes desagregado en sus siete componentes, para el nivel
 de tensión 2 en categoría oficial. A la derecha, el peso medio de cada
 componente sobre el horizonte.
 \fuente{elaboración propia desde
 \file{data/tarifas\_cedenar\_mensual.csv}, transcrito de los boletines
 mensuales del comercializador archivados en
 \file{data/cedenar\_pdfs/}.}}
 \label{fig:precios-cu}
\end{figure}

\notafig{El costo unitario medio del horizonte es de
\uni{792.06}{COP/kWh} en categoría oficial, con un recorrido anual de
\uni{50.18}{COP/kWh}, mínimo en enero de 2026 y máximo en agosto de 2025.
La estabilidad de la barra apilada a lo largo de los trece meses indica
que ninguna componente sufre saltos regulatorios dentro del horizonte.
Elaboración propia desde \file{data/tarifas\_cedenar\_mensual.csv}.}

Su reparto medio es:
generación, \uni{305.6}{COP/kWh} (\pct{38.6}); comercialización,
\uni{174.6}{} (\pct{22.0}); distribución, \uni{157.6}{} (\pct{19.9});
transmisión, \uni{54.9}{} (\pct{6.9}); otros cargos, \uni{41.2}{}
(\pct{5.2}); restricciones, \uni{36.8}{} (\pct{4.6}); y pérdidas,
\uni{21.3}{} (\pct{2.7}).

\begin{lecturabox}
La cifra que conviene retener de esta figura es que \textbf{la generación
apenas explica el \pct{38.6} de lo que se paga}. El resto, casi dos
tercios, son cargos de red, comercialización y conceptos regulatorios. Ahí
está la razón económica de fondo de todo el trabajo: un intercambio que
ocurre dentro de la comunidad, sin usar la red de distribución, puede
evitar una fracción de esos dos tercios, y esa fracción evitada es mayor
que cualquier ganancia de eficiencia en la generación misma. El
Capítulo~\ref{sec:disc} vuelve sobre este punto porque es también la
objeción más seria que se le puede hacer al mecanismo.
\end{lecturabox}

\subsection{Dos tarifas para la misma comunidad}
\label{sub:precios-categorias}

Las cinco instituciones no pagan lo mismo. Dos de ellas, Udenar y el
HUDN, son entidades públicas y pagan el costo unitario exacto; las otras
tres (Mariana, la UCC y Cesmag) son privadas y pagan además la
contribución de solidaridad del veinte por ciento. La
Figura~\ref{fig:precios-categorias} pone las dos series sobre un mismo eje
y sombrea la distancia entre ellas.

\begin{figure}[H]
 \centering
 \includegraphics[width=0.95\textwidth]{figuras/f5_04_oficial_comercial.png}
 \caption[Costo unitario por categoría tarifaria]{Costo unitario
 mensual de las dos categorías presentes en la comunidad. La franja
 sombreada entre ambas es la contribución de solidaridad.
 \fuente{elaboración propia desde
 \file{data/tarifas\_cedenar\_mensual.csv}.}}
 \label{fig:precios-categorias}
\end{figure}

\notafig{La diferencia media es de \uni{158.4}{COP/kWh}, exactamente el
\pct{20.0} del costo oficial. La razón entre las dos series vale
$1{,}2000$ en los trece meses, sin una sola excepción, lo que confirma que
se trata de un factor regulatorio y no de una diferencia de condiciones de
suministro. Elaboración propia desde
\file{data/tarifas\_cedenar\_mensual.csv}.}

Esta heterogeneidad tiene una consecuencia de diseño que conviene dejar
explícita, porque es una de las decisiones metodológicas menos evidentes
del trabajo: en la liquidación de los escenarios regulatorios, el precio
al que cada institución compra no es un número sino una matriz, con una
fila por institución y una columna por hora, que asigna a cada hora el
costo unitario vigente para esa institución en el mes que la contiene. Una
comunidad de cinco miembros con dos categorías tarifarias y trece meses de
tarifas distintas no puede liquidarse con un escalar sin introducir un
sesgo sistemático.

\begin{detallebox}
El juego del Capítulo~\ref{sec:modelo}, en cambio, sí usa un escalar
comunitario. La asimetría es deliberada y conviene entenderla: la
dinámica que busca el precio de equilibrio necesita un único techo común
(si cada agente tuviera el suyo, no habría un precio de mercado sino
cinco), mientras que la liquidación posterior sí puede y debe respetar la
tarifa de cada uno. El mecanismo que reconcilia ambas cosas es un techo
por agente aplicado en el momento de liquidar, que el
Capítulo~\ref{sec:modelo} desarrolla.
\end{detallebox}

\subsection{La banda de ganancia mutua}
\label{sub:precios-banda}

Con las dos series construidas queda delimitado el espacio dentro del cual
un intercambio entre pares tiene sentido económico. La condición es doble
y elemental: el precio debe superar lo que el vendedor obtendría
inyectando a la red, y quedar por debajo de lo que el comprador pagaría
comprándola a su comercializador. Si se cumple, las dos partes ganan; si
no se cumple, al menos una prefiere la red. La
Figura~\ref{fig:precios-banda} reúne los cuatro niveles construidos en
este capítulo y delimita con ellos la banda resultante.

\begin{figure}[H]
 \centering
 \includegraphics[width=0.95\textwidth]{figuras/f5_07_banda.png}
 \caption[Piso y techo del precio entre pares]{Los cuatro niveles de
 referencia construidos en este capítulo y la banda que delimitan. Todo
 precio de equilibrio del mercado entre pares vive dentro de la franja
 sombreada.
 \fuente{elaboración propia desde
 \file{data/precios\_bolsa\_xm\_api.csv} y
 \file{data/tarifas\_cedenar\_mensual.csv}.}}
 \label{fig:precios-banda}
\end{figure}

\notafig{La banda no procede de una diferencia de costos de producción
sino de la distancia entre lo que la red paga por recibir energía y lo que
cobra por entregarla. Es, por tanto, un margen de origen regulatorio y no
tecnológico, lo que condiciona su permanencia: una reforma tarifaria puede
estrecharlo sin que cambie nada en las instalaciones. Elaboración propia
desde las series de bolsa y de tarifa.}

La amplitud de esa banda es lo que hace viable el ejercicio, y su origen
es enteramente regulatorio: no procede de una diferencia de costos de
producción sino de la distancia entre lo que la red paga por recibir
energía y lo que cobra por entregarla. Entre la media de bolsa
(\uni{182.5}{COP/kWh}) y el costo unitario oficial, \uni{792.06}{}, hay un
factor superior a cuatro. Cualquier precio intermedio deja a las dos
partes mejor que la alternativa de operar contra la red.

\begin{trampabox}
A lo largo del documento aparecen varias cifras asociadas al «precio al
que la red compra» y conviene deshacer la confusión de una vez, porque
corresponden a papeles distintos de un mismo referente económico. El
\emph{piso de la dinámica} del juego es una constante de
\uni{280}{COP/kWh}, heredada de la calibración inicial, que acota
inferiormente la búsqueda del equilibrio. La \emph{serie horaria de
bolsa}, con media de \uni{182.5}{COP/kWh}, no interviene en esa dinámica:
es la que liquida los escenarios que la requieren y la exportación a red
del excedente no vendido. Y el \emph{valor nominal} de los barridos de
sensibilidad del Capítulo~\ref{sec:rob} es también \uni{280}{}, coherente
con el piso del juego. Piso de la dinámica, serie de liquidación y nominal
de barridos son tres papeles y no deben confundirse.
\end{trampabox}

\subsection{Lo que cuesta generar}
\label{sub:precios-lcoe}

Falta un tercer precio, que no se compra ni se vende: lo que le cuesta a
la propia institución el kilovatio-hora que produce. Para un sistema
fotovoltaico ya instalado, el costo marginal de producir una unidad
adicional es prácticamente nulo, de modo que el parámetro relevante es el
costo nivelado, es decir, la inversión amortizada sobre la energía que el
sistema entregará a lo largo de su vida.

El valor adoptado es de \uni{225}{COP/kWh} para las cuatro instituciones
con inversor de un fabricante y \uni{210}{} para Cesmag, que tiene otro.
Tras el ajuste por la irradiancia local de Pasto, los valores efectivos en
el simulador quedan en torno a \uni{241}{COP/kWh} para las primeras y
\uni{225}{} para Cesmag, lo que introduce una heterogeneidad del
\pct{6.67} entre ambas.

El rango se sustenta en tres referencias de distinto alcance: el costo
nivelado ponderado global de la energía solar fotovoltaica que reporta la
agencia internacional de energías renovables; los costos de referencia
para generación solar en Colombia del plan indicativo de expansión, entre
\num{300} y \uni{450}{COP/kWh} para escalas utilitaria y media; y el
análisis de incentivos económicos a la energía solar en el país, que sitúa
el costo nivelado efectivo de sistemas distribuidos en el rango de
\num{200} a \uni{250}{COP/kWh} bajo condiciones de financiación
favorables.

\begin{detallebox}
El valor adoptado es conservador dentro de ese rango, y esa elección tiene
una consecuencia contraintuitiva que el Capítulo~\ref{sec:rob} demuestra
de forma cuantitativa: el resultado del estudio \textbf{no depende de este
parámetro}. El análisis de sensibilidad global asigna al costo de generar
un índice de influencia sobre la ganancia del orden de $10^{-7}$, es decir,
indistinguible de cero. La razón es estructural: lo que el mercado transa
es excedente después de autoconsumo, y el volumen de ese excedente lo
determina el recurso disponible, no el costo contable de haberlo
producido. Se documenta aquí, no obstante, con el mismo detalle que el
resto, porque un parámetro que resulta irrelevante \emph{a posteriori} no
podía saberse irrelevante \emph{a priori}.
\end{detallebox}

\subsection{Una provisionalidad declarada}
\label{sub:precios-provisional}

El nivel de tensión asumido para las cinco instituciones, del cual
dependen todas las tarifas de este capítulo, es el nivel 2, y es
provisional hasta su verificación contra facturas reales. La publicación
tarifaria del comercializador respalda los valores de la categoría no
residencial en ese nivel y el factor de $1{,}20$ entre las dos categorías;
lo que queda abierto es únicamente la verificación del nivel de tensión
efectivo de cada institución, que la revisión regulatoria del proyecto
sitúa en el nivel 1.

Se declara aquí y no en un anexo porque afecta al valor absoluto de todas
las cifras económicas del documento. No afecta, en cambio, al
ordenamiento entre mecanismos, que es lo que el trabajo pretende
establecer: un cambio de nivel de tensión desplaza el techo de la banda
para las cinco instituciones por igual, y el Capítulo~\ref{sec:rob}
muestra que el ordenamiento sobrevive a variaciones de ese techo muy
superiores a la diferencia entre los dos niveles.

En síntesis, con este capítulo queda completo el insumo del modelo: a las
dos matrices de energía del Capítulo~\ref{sec:prep} se suman ahora la
serie horaria de bolsa, la matriz mensual de tarifa por institución y el
costo nivelado de generar, junto con la banda de \uni{182.5}{} a
\uni{792.06}{COP/kWh} dentro de la cual puede existir un intercambio que
convenga a las dos partes. El capítulo siguiente describe el mecanismo que
busca un precio dentro de esa banda.


% =====================================================================
%  PARTE II — EL INSTRUMENTO
% =====================================================================
% =====================================================================
%  El modelo P2P
%  Parte II — El instrumento
%
%  ENCARGO: Formulación del juego de Stackelberg con dinámica de replicador y relajación lagrangiana, pieza por pieza.
%
%  Reglas del documento:
%   - Una idea = una subseccion = una figura.
%   - Toda figura declara su procedencia con \fuente{...}.
%   - Toda cifra sale del canon de agosto; ver GUIA.md.
%   - M1 y M3 se reportan siempre en paralelo.
%   - Ningun superlativo sin releer el CSV de la frontera correcta.
% =====================================================================
\section{El modelo P2P}
\label{sec:modelo}
\justifying

\placeholderfig{Capítulo pendiente de redacción. Formulación del juego de Stackelberg con dinámica de replicador y relajación lagrangiana, pieza por pieza.}

% =====================================================================
%  Verificación del solucionador
%  Parte II — El instrumento
%
%  ENCARGO: Prueba dorada contra el oráculo, contraste MATLAB frente a Python, el ciclo de período dos y la sensibilidad al horizonte de integración.
%
%  Reglas del documento:
%   - Una idea = una subseccion = una figura.
%   - Toda figura declara su procedencia con.
%   - Toda cifra sale del canon de agosto; ver GUIA.md.
%   - M1 y M3 se reportan siempre en paralelo.
%   - Ningun superlativo sin releer el CSV de la frontera correcta.
% =====================================================================
\section{Verificación del solucionador}
\label{sec:verif}
\justifying

\placeholderfig{Capítulo pendiente de redacción. Prueba dorada contra el oráculo, contraste MATLAB frente a Python, el ciclo de período dos y la sensibilidad al horizonte de integración.}

% =====================================================================
%  Los escenarios regulatorios como algoritmos
%  Parte II — El instrumento
%
%  ENCARGO: C1 a C5 y el P2P como sexto mecanismo. Cada uno con su norma, su pseudocódigo horario y su diagrama de flujo.
%
%  Reglas del documento:
%   - Una idea = una subseccion = una figura.
%   - Toda figura declara su procedencia con.
%   - Toda cifra sale del canon de agosto; ver GUIA.md.
%   - M1 y M3 se reportan siempre en paralelo.
%   - Ningun superlativo sin releer el CSV de la frontera correcta.
% =====================================================================
\section{Los escenarios regulatorios como algoritmos}
\label{sec:esc}
\justifying

\placeholderfig{Capítulo pendiente de redacción. C1 a C5 y el P2P como sexto mecanismo. Cada uno con su norma, su pseudocódigo horario y su diagrama de flujo.}


% =====================================================================
%  PARTE III — LA DEPURACIÓN
% =====================================================================
% =====================================================================
%  Cómo evolucionó el modelo
%  Parte III — La depuración
%
%  ENCARGO: Las decisiones que movieron las cifras, cada una con su figura de impacto, y los hallazgos que se cayeron al comprobarlos.
%
%  Reglas del documento:
%   - Una idea = una subseccion = una figura.
%   - Toda figura declara su procedencia con.
%   - Toda cifra sale del canon de agosto; ver GUIA.md.
%   - M1 y M3 se reportan siempre en paralelo.
%   - Ningun superlativo sin releer el CSV de la frontera correcta.
% =====================================================================
\section{Cómo evolucionó el modelo}
\label{sec:depu}
\justifying

\placeholderfig{Capítulo pendiente de redacción. Las decisiones que movieron las cifras, cada una con su figura de impacto, y los hallazgos que se cayeron al comprobarlos.}


% =====================================================================
%  PARTE IV — LA EVIDENCIA
% =====================================================================
% =====================================================================
%  Resultados: del flujo horario a la cifra agregada
%  Parte IV — La evidencia
%
%  ENCARGO: Orden estrictamente ascendente: hora, día, mes, agregado. Nunca al revés.
%
%  Reglas del documento:
%   - Una idea = una subseccion = una figura.
%   - Toda figura declara su procedencia con.
%   - Toda cifra sale del canon de agosto; ver GUIA.md.
%   - M1 y M3 se reportan siempre en paralelo.
%   - Ningun superlativo sin releer el CSV de la frontera correcta.
% =====================================================================
\section{Resultados: del flujo horario a la cifra agregada}
\label{sec:res}
\justifying

\placeholderfig{Capítulo pendiente de redacción. Orden estrictamente ascendente: hora, día, mes, agregado. Nunca al revés.}

% =====================================================================
%  Robustez: sensibilidad local, bivariada y global
%  Parte IV — La evidencia
%
%  ENCARGO: Barridos uniparamétricos, mapa bivariado e índices de Sobol sobre el caso de estudio real.
%
%  Reglas del documento:
%   - Una idea = una subseccion = una figura.
%   - Toda figura declara su procedencia con \fuente{...}.
%   - Toda cifra sale del canon de agosto; ver GUIA.md.
%   - M1 y M3 se reportan siempre en paralelo.
%   - Ningun superlativo sin releer el CSV de la frontera correcta.
% =====================================================================
\section{Robustez: sensibilidad local, bivariada y global}
\label{sec:rob}
\justifying

\placeholderfig{Capítulo pendiente de redacción. Barridos uniparamétricos, mapa bivariado e índices de Sobol sobre el caso de estudio real.}

% =====================================================================
%  Factibilidad, racionalidad individual y equidad
%  Parte IV — La evidencia
%
%  ENCARGO: Cumplimiento regulatorio, umbral de deserción por agente, robustez ante retiro y precio de la equidad.
%
%  Reglas del documento:
%   - Una idea = una subseccion = una figura.
%   - Toda figura declara su procedencia con.
%   - Toda cifra sale del canon de agosto; ver GUIA.md.
%   - M1 y M3 se reportan siempre en paralelo.
%   - Ningun superlativo sin releer el CSV de la frontera correcta.
% =====================================================================
\section{Factibilidad, racionalidad individual y equidad}
\label{sec:fact}
\justifying

\placeholderfig{Capítulo pendiente de redacción. Cumplimiento regulatorio, umbral de deserción por agente, robustez ante retiro y precio de la equidad.}

% =====================================================================
%  Inferencia estadística
%  Parte IV — La evidencia
%
%  ENCARGO: Remuestreo por bloques sobre la serie diaria, intervalo de confianza y tamaño del efecto.
%
%  Reglas del documento:
%   - Una idea = una subseccion = una figura.
%   - Toda figura declara su procedencia con.
%   - Toda cifra sale del canon de agosto; ver GUIA.md.
%   - M1 y M3 se reportan siempre en paralelo.
%   - Ningun superlativo sin releer el CSV de la frontera correcta.
% =====================================================================
\section{Inferencia estadística}
\label{sec:inf}
\justifying

\placeholderfig{Capítulo pendiente de redacción. Remuestreo por bloques sobre la serie diaria, intervalo de confianza y tamaño del efecto.}


% =====================================================================
%  PARTE V — CIERRE
% =====================================================================
% =====================================================================
%  Discusión
%  Parte V — Cierre
%
%  ENCARGO: Lectura económica y regulatoria de lo obtenido, limitaciones y recomendaciones de política.
%
%  Reglas del documento:
%   - Una idea = una subseccion = una figura.
%   - Toda figura declara su procedencia con.
%   - Toda cifra sale del canon de agosto; ver GUIA.md.
%   - M1 y M3 se reportan siempre en paralelo.
%   - Ningun superlativo sin releer el CSV de la frontera correcta.
% =====================================================================
\section{Discusión}
\label{sec:disc}
\justifying

\placeholderfig{Capítulo pendiente de redacción. Lectura económica y regulatoria de lo obtenido, limitaciones y recomendaciones de política.}

% =====================================================================
%  Conclusiones y trabajo futuro
%  Parte V — Cierre
%
%  ENCARGO: En espejo con los objetivos específicos, un cierre por objetivo.
%
%  Reglas del documento:
%   - Una idea = una subseccion = una figura.
%   - Toda figura declara su procedencia con \fuente{...}.
%   - Toda cifra sale del canon de agosto; ver GUIA.md.
%   - M1 y M3 se reportan siempre en paralelo.
%   - Ningun superlativo sin releer el CSV de la frontera correcta.
% =====================================================================
\section{Conclusiones y trabajo futuro}
\label{sec:conc}
\justifying

\placeholderfig{Capítulo pendiente de redacción. En espejo con los objetivos específicos, un cierre por objetivo.}


% =====================================================================
%  BIBLIOGRAFÍA
% =====================================================================
\clearpage
\phantomsection
\addcontentsline{toc}{section}{Referencias}
\bibliographystyle{IEEEtran}
\bibliography{referencias}

% =====================================================================
%  ANEXOS
% =====================================================================
\iniciaranexos
% =====================================================================
%  Reproducibilidad de la corrida canónica
%  Anexos
%
%  ENCARGO: Comandos exactos, entorno congelado, huellas de los artefactos y compuertas de verificación.
%
%  Reglas del documento:
%   - Una idea = una subseccion = una figura.
%   - Toda figura declara su procedencia con \fuente{...}.
%   - Toda cifra sale del canon de agosto; ver GUIA.md.
%   - M1 y M3 se reportan siempre en paralelo.
%   - Ningun superlativo sin releer el CSV de la frontera correcta.
% =====================================================================
\section{Reproducibilidad de la corrida canónica}
\label{anx:repro}
\justifying

\subsection{Cómo se localizan y se leen los archivos crudos}
\label{sub:repro-lectura}

El Capítulo~\ref{sec:prep} resume en una sola etapa lo que el cargador
hace para llegar a la serie de cada equipo. Se llama así aquí a la pieza
que recorre los archivos, los lee y entrega la serie horaria de un equipo. El detalle queda aquí, porque
describe el programa y no el dato. De las decisiones que siguen, casi
ninguna llega a ejercerse sobre este conjunto, y la única que sí se
ejerce no cambia ningún valor.

\textbf{La búsqueda de carpetas tolera diferencias de escritura.} Los
nombres con que llegan las carpetas no son uniformes, y el de los
medidores de la Universidad CESMAG viene con una errata de origen, escrita
\code{eletricMeter}. Se conserva tal cual, para reproducir la estructura de
la fuente en vez de corregirla y perder la trazabilidad. La tolerancia, sin
embargo, no llega a hacer falta: las diez carpetas que el modelo emplea
coinciden exactamente con lo declarado en la configuración.

\textbf{Cada archivo se decodifica probando cuatro codificaciones en
orden}: \code{utf-8-sig}, \code{utf-8}, \code{cp1252} y \code{latin-1}. El
orden importa porque \code{latin-1} acepta cualquier secuencia de bytes y
nunca falla, de modo que probarla antes que \code{cp1252} la haría ganar
siempre y corrompería los caracteres del rango \code{0x80}--\code{0x9F},
que es donde viven las tildes.

Tampoco esa cascada se ejerce. Los 73 archivos del árbol resuelven en el
primer candidato, y ninguno contiene un solo byte por encima de 127. Se
documenta porque el orden es una decisión deliberada del cargador y porque
un conjunto de datos posterior sí podría ejercerla, no porque haya
cambiado una cifra de este trabajo.

\textbf{La fecha y el valor se convierten en modo tolerante.} Lo que no se
deja interpretar queda nulo en vez de detener la carga, y las dos columnas
no reciben el mismo trato: una lectura sin fecha válida se descarta entera,
porque no tiene dónde colocarse, mientras que una con fecha válida y valor
ilegible se conserva en blanco para que la limpieza la rellene más
adelante. Ninguna de las dos conversiones llega a descartar nada, porque
sobre los \num{2813305} registros de los diez medidores que el modelo
emplea no hay una sola fecha ilegible ni un solo valor no numérico.

\textbf{Los instantes repetidos se funden promediando.} La fusión ocurre
dentro de cada archivo, antes de alinear los tramos de un mismo medidor, y
su motivo es de forma y no de contenido: el índice ha de ser único para que
los tramos se puedan superponer sobre un mismo eje de tiempo. La
Tabla~\ref{tab:repro-dup} recoge dónde ocurre.

\begin{table}[H]
 \centering
 \caption{Instantes con más de una lectura, por institución y frontera. Las
 instituciones que no aparecen no tienen ninguno.}
 \label{tab:repro-dup}
 \begin{tabular}{@{}l l r r@{}}
 \toprule
 \textbf{Institución} & \textbf{Frontera} &
 \textbf{Instantes repetidos} & \textbf{Máximo en un instante} \\
 \midrule
 UCC & M1 & \num{3071} & 11 \\
 UCC & M3 & \num{3551} & 73 \\
 \addlinespace
 Mariana & M1 & \num{28} & 2 \\
 Mariana & M3 & \num{71} & 4 \\
 \bottomrule
 \end{tabular}
\end{table}

Las tres reglas de fusión posibles, promediar, quedarse con la primera o
quedarse con la última, dan el mismo resultado sobre este conjunto, porque
las lecturas que comparten instante coinciden hasta el último dígito en la
totalidad de los casos. Lo que sí informa es la concentración: la UCC
aporta el \pct{98.5} de los instantes repetidos, y su medidor secundario
llega a volcar 73 lecturas idénticas en un mismo instante. Es un rasgo del
equipo y no del método, y no altera ninguna cifra.

\textbf{La unidad se convierte al final, y ahí el orden solo importa para
el recorte.} El cargador promedia a paso horario y solo después divide
entre mil la serie del inversor. Para la división el orden es indiferente,
porque dividir entre una constante y promediar se pueden intercambiar sin
que el resultado cambie; comprobado sobre los siete inversores, invertirlo
produce diferencias a partir de la decimoquinta cifra decimal, que es ruido
de redondeo. Con el recorte a valores no negativos no ocurre lo mismo:
recortar y promediar no se pueden intercambiar, porque una lectura negativa
dentro de una hora de media positiva sobreviviría al promedio y rebajaría
el valor horario. El cargador recorta sobre la serie ya horaria, de modo
que ese caso podría darse; sobre este conjunto no llega a existir, porque
ninguna de las \num{793981} lecturas crudas de los siete inversores es
negativa.

\textbf{Los tres tipos de medidor llevan otro nombre en el canon.} El
Capítulo~\ref{sec:prep} los llama neto, neto parcial y bruto, que es el
vocabulario del texto y de las figuras. En la configuración del cargador y
en las tablas de datos que acompañan a cada figura figuran como
\code{net}, \code{net\_partial} y \code{gross}. Se anota la
correspondencia aquí para que quien audite el código no la tenga que
adivinar.

\textbf{El orden exacto} en que se aplica todo lo anterior es el siguiente:
localizar la carpeta del equipo, decodificar cada archivo, convertir la
fecha y el valor, fundir los instantes repetidos, alinear los tramos del
medidor sobre un mismo eje de tiempo, recortar al horizonte, promediar a
paso horario y, en último lugar, convertir la unidad.

\textbf{Cómo se lee la figura del recorrido.} El
Capítulo~\ref{sec:prep} traza una hora real desde el archivo hasta la serie
horaria, y conviene decir qué hay en cada banda de la
Figura~\ref{fig:prep-archivo-serie}. Arriba, los archivos del medidor y del
inversor sobre el calendario, con el horizonte de estudio sombreado. En el
centro, las lecturas a paso de dos minutos de esa hora: el medidor entrega
38 filas que al depurar quedan en 30 instantes, ocho de los cuales traen la
lectura duplicada, y el inversor entrega 30 lecturas en vatios, ninguna
negativa, de modo que el recorte a cero no llega a actuar. Abajo, la media
de cada panel ocupa su lugar entre las 24 horas del día.

Fundir los duplicados quita una copia y no rellena nada. Los instantes que
quedan caen en los treinta tramos de dos minutos de la hora sin faltar
ninguno, de modo que el ejemplo ilustra el caso en que la medida cubre el
intervalo entero y no hay tramo que suponer.

\subsection{El criterio del regulador para la lectura que falta}
\label{sub:repro-norma}

El Capítulo~\ref{sec:prep} invoca la norma como criterio y no como
obligación. Aquí queda el articulado que lo sostiene, porque la analogía
solo vale si se puede comprobar.

El artículo 38 de la Resolución CREG 038 de 2014 se activa mientras se
reparan o reponen los elementos de un sistema de medición en falla o
hurtado, y ordena estimar la lectura por dos vías distintas según la
frontera. Si la frontera reporta al administrador del sistema de
intercambios comerciales, integrando la medida de potencia activa cuando
esa medida está bajo la supervisión de un centro de control, o recurriendo
a curvas típicas. Si no reporta, conforme a la Ley 142 de 1994 y al
artículo 31 de la Resolución CREG 108 de 1997, que ordena establecer el
consumo con base en los consumos promedios de otros períodos del mismo
usuario, en los de usuarios en circunstancias similares o en aforos
individuales.

Dos precisiones importan. La primera es que \textbf{ninguno de los medios
que el artículo enumera consiste en asignar cero}, y conviene decirlo así
porque es una afirmación sobre la lista y no sobre una prohibición que la
norma no formula. La segunda es que la vía aplicable por analogía a este
trabajo es la segunda y no la primera: integrar la potencia activa significa
allí leer otro canal que sí funciona, mientras que estimar por promedios de
otros períodos del mismo usuario es exactamente lo que aquí se hace con la
hora incompleta.

Nada de esto obliga. Los equipos del proyecto son instrumentación de
investigación y no fronteras comerciales, y lo que interrumpe el registro no
es un sistema de medición en falla sino un corte de telemetría.

El texto del artículo 38 se verificó por dos vías independientes que
coinciden palabra por palabra. Se deja anotado porque la publicación oficial
en línea de la resolución se trunca dentro del artículo 28 y no alcanza al
38, de modo que quien quiera comprobarlo no lo encontrará en la fuente
evidente.

\subsection{Cómo se lee la figura de la profundidad de la lectura}
\label{sub:repro-profundidad}

La Figura~\ref{fig:prep-gradacion} del Capítulo~\ref{sec:prep} resume cada
uno de los cinco medidores de la frontera principal en un renglón, y la
distancia entre sus dos marcas es la gradación que la subsección describe.
Fila por fila:

En Udenar el medidor baja bastante más abajo del cero de lo que marca de
ordinario, es decir, la inversión del flujo no es allí un episodio sino el
estado normal de las horas de sol.

En Mariana y en la UCC el descenso no llega siquiera a igualar la lectura
mediana de su propio circuito, de manera que el flujo se invierte solo
cuando la producción culmina y la carga afloja.

En el HUDN el recorrido entero se queda por encima de \uni{6}{kW}, y en
CESMAG termina en \uni{0.195}{kW}, es decir, roza el cero sin cruzarlo,
porque detrás de esos dos medidores no genera nada.

\placeholderfig{Anexo pendiente de redacción. Comandos exactos, entorno congelado, huellas de los artefactos y compuertas de verificación.}

% =====================================================================
%  Registro completo de decisiones técnicas
%  Anexos
%
%  ENCARGO: Las decisiones fechadas CAL-1 a CAL-43, con su fecha, su motivo y su efecto.
%
%  Reglas del documento:
%   - Una idea = una subseccion = una figura.
%   - Toda figura declara su procedencia con \fuente{...}.
%   - Toda cifra sale del canon de agosto; ver GUIA.md.
%   - M1 y M3 se reportan siempre en paralelo.
%   - Ningun superlativo sin releer el CSV de la frontera correcta.
% =====================================================================
\section{Registro completo de decisiones técnicas}
\label{anx:cal}
\justifying

\placeholderfig{Anexo pendiente de redacción. Las decisiones fechadas CAL-1 a CAL-43, con su fecha, su motivo y su efecto.}

% =====================================================================
%  Trazabilidad con las actividades de la propuesta
%  Anexos
%
%  ENCARGO: Correspondencia entre cada actividad de la propuesta aprobada y la sección que la sustenta.
%
%  Reglas del documento:
%   - Una idea = una subseccion = una figura.
%   - Toda figura declara su procedencia con.
%   - Toda cifra sale del canon de agosto; ver GUIA.md.
%   - M1 y M3 se reportan siempre en paralelo.
%   - Ningun superlativo sin releer el CSV de la frontera correcta.
% =====================================================================
\section{Trazabilidad con las actividades de la propuesta}
\label{anx:traza}
\justifying

\placeholderfig{Anexo pendiente de redacción. Correspondencia entre cada actividad de la propuesta aprobada y la sección que la sustenta.}

% =====================================================================
%  Inventario de datos y artefactos
%  Anexos
%
%  ENCARGO: Qué archivo es cada cosa, cuál es canónico y cuál está retirado.
%
%  Reglas del documento:
%   - Una idea = una subseccion = una figura.
%   - Toda figura declara su procedencia con \fuente{...}.
%   - Toda cifra sale del canon de agosto; ver GUIA.md.
%   - M1 y M3 se reportan siempre en paralelo.
%   - Ningun superlativo sin releer el CSV de la frontera correcta.
% =====================================================================
\section{Inventario de datos y artefactos}
\label{anx:inv}
\justifying

\placeholderfig{Anexo pendiente de redacción. Qué archivo es cada cosa, cuál es canónico y cuál está retirado.}

% =====================================================================
%  Preguntas abiertas al comité
%  Anexos
%
%  ENCARGO: Las decisiones que dependen de los asesores y su estado.
%
%  Reglas del documento:
%   - Una idea = una subseccion = una figura.
%   - Toda figura declara su procedencia con.
%   - Toda cifra sale del canon de agosto; ver GUIA.md.
%   - M1 y M3 se reportan siempre en paralelo.
%   - Ningun superlativo sin releer el CSV de la frontera correcta.
% =====================================================================
\section{Preguntas abiertas al comité}
\label{anx:preg}
\justifying

Este anexo reúne las decisiones que no corresponden a este trabajo y que se
dejan planteadas con la evidencia reunida, para que quien decida lo haga
sobre cifras y no sobre impresiones. Cada pregunta se enuncia, se acompaña
de lo medido y termina con lo que cambiaría según la respuesta.

\subsection{Las fuentes del modelo base y cuál manda}
\label{sub:preg-competencia}

Esta entrada nació como pregunta y dejó de serlo al aparecer la versión
arbitrada del modelo. Se conserva porque el desacuerdo entre fuentes sigue
existiendo y conviene que quede constancia de cuál se siguió.

El bienestar del comprador incluye un término que penaliza competir con los
demás, y un segundo término que recoge lo que paga por la energía. El guion
en Python que acompaña al modelo escribe el primero multiplicando por la
energía del comprador propio, mientras que la versión arbitrada, publicada
en \emph{IEEE Latin America Transactions} en agosto de 2025, lo escribe
multiplicando por la energía del comprador ajeno, igual que la formulación
que aparece comentada en el código fuente en Matlab. Para el segundo, el
documento extenso del modelo introduce un factor del precio de compra a la
red que ni la versión arbitrada ni el guion contemplan.

Las dos expresiones del primer término coinciden exactamente cuando todos
los compradores adquieren la misma cantidad, que es el caso de las
verificaciones sintéticas con que se validó la traducción en su momento, y
se separan en cuanto las cantidades difieren. La magnitud de la separación
se midió sobre cuatro compradores con cantidades distintas: la diferencia
mayor asciende a \num{1755,95} unidades sobre valores del orden de
\num{2500}.

\textbf{Este trabajo sigue en los dos casos a la versión arbitrada}, que es
la especificación publicada. La elección está además aislada: se comprobó
que ninguna de las funciones que resuelven el mercado consulta esas
expresiones, de modo que no mueve ni un flujo ni un precio. De los \num{672}
valores que componen la huella de verificación de una jornada completa,
cambiar de forma altera \num{25}, y los veinticinco son bienestar del
comprador.

Queda por confirmar, y es lo único que sigue abierto en este punto, si el
guion en Python recoge una versión anterior del modelo o una corrección
posterior que no llegó al artículo.

\subsection{Las dos preguntas regulatorias}
\label{sub:preg-regulatorias}

Quedan planteadas y sin resolver desde el repositorio.

La primera es si la exención de la contribución de solidaridad está
efectivamente radicada para las cinco instituciones, o si es solo aplicable
en derecho. El trabajo la aplica a las cinco, por tratarse de usuarios no
regulados y de servicios de salud y educación, pero eso es una lectura de la
norma y no la verificación de un trámite.

La segunda es si la energía intercambiada dentro de la comunidad paga o no
cargos por uso de redes. La respuesta gobierna el ancho de la banda de
precios y, con él, el valor entero del mercado: sin exención la banda se
vacía en la gran mayoría de las horas.

Para esta segunda conviene llevar un antecedente del propio grupo. El
trabajo sobre mercados entre pares en escenarios multimicrorred, publicado
en \emph{TecnoLógicas} en 2024, modela explícitamente una penalización, o
costo de transmisión, cuando un recurso envía energía a una microrred
vecina. Es decir, el grupo ya ha modelado que el transporte entre zonas se
cobra, lo que sitúa la pregunta en un terreno conocido en lugar de abrirla
desde cero.

\subsection{Si el techo del precio es uno solo o uno por comprador}
\label{sub:preg-techo}

El límite superior del precio del mercado es el costo unitario que cada
institución le paga a su comercializador, porque por encima de él le
conviene comprarle a la red. La comunidad tiene dos comercializadores, de
modo que ese límite no es el mismo para las cinco: cuatro instituciones lo
tienen en 731,1 (COP/kWh) y la quinta en 777,2.

La consulta al asesor externo en materia regulatoria, de junio de 2026,
parece resolver el asunto y no lo resuelve. Su respuesta afirma tres veces
que el costo unitario base es el mismo para todos, pero contesta por clase
de usuario, es decir comercial frente a oficial dentro de un mismo
comercializador, y esa distinción ya se retiró al reconocer a las cinco como
usuarios no regulados. Lo que hoy separa los dos techos es el
comercializador, y sobre ese eje no se preguntó.

Se midieron los dos regímenes sobre sesenta horas activas de cada frontera,
evaluando siempre el ahorro contra el costo unitario real de cada
institución. La energía transada es idéntica en ambos, con diferencias del
orden de una diezbillonésima de kilovatio hora, lo que confirma que el
límite superior no interviene en la cantidad. Lo que cambia es el
emparejamiento y el reparto.

Con el techo de cada comprador, el excedente de la frontera principal
asciende a 46.596,6 (COP) y el 82 \% del ahorro del lado comprador se
concentra en la institución del otro comercializador, porque la energía
escasa va hacia quien tiene la alternativa externa más cara. Con un techo
único e igual al más bajo, el excedente se sitúa en 45.016,1 (COP), un
3,4 \% menos, y el ahorro se reparte de forma más pareja. La disyuntiva, por
tanto, es entre eficiencia y equidad, y no entre correcto e incorrecto.

Un tercer régimen queda descartado sin necesidad de decisión: un techo único
e igual al más alto obliga a las cuatro instituciones del techo bajo a pagar
756,1 (COP/kWh) por energía que la red les vendería a 731,1, con lo que su
ahorro se vuelve negativo en las ochenta y ocho horas medidas. Ningún
régimen que obligue a un participante a pagar por encima de su alternativa
externa es admisible.

La implementación admite las dos respuestas. Cuando los techos coinciden, la
formulación por comprador se reduce exactamente a la de techo único, de
manera que la decisión del comité no obliga a reescribir nada.

\subsection{Cómo se liquida un mercado horario con un crédito mensual}
\label{sub:preg-liquidacion}

Esta entrada nació como pregunta y la resolvió el propio texto regulatorio. El
Anexo~4 de la Resolución CREG 101 072 de 2025, que aplica el artículo 26 de la
Resolución CREG 174 durante el periodo de transición, liquida el crédito con
las cantidades del mes y valora a la bolsa de cada hora lo que lo supera, a
partir de una hora de corte que define como aquella en que los excedentes
horarios acumulados desde el inicio del mes igualan o sobrepasan la
importación total del mes. El cupo se fija con el total mensual, la hora del
corte con el acumulado horario, y el dinero se liquida al cierre del período.

Un mercado entre pares encaja en esa misma estructura: los intercambios se
registran por hora, se liquidan al cierre del mes junto con la factura, y el
piso de cada vendedor es la permuta antes de su hora de corte y el precio de
bolsa de la hora a partir de ella. Como la hora de corte depende de la
importación total, solo se conoce al cerrar el mes; en una liquidación hecha
sobre datos medidos eso equivale a suponer conocidos los totales del mes, y
así se declara. En la frontera principal ningún mes alcanza la hora de corte;
en la secundaria, el reparto del exceso entre las horas apenas altera su
valor.

\subsection{El precio del excedente que supera el crédito}
\label{sub:preg-mcm}

Tampoco queda abierta. La regla definitiva de la Resolución CREG 101 072 de
2025 valora ese excedente a la variable MCm, el costo promedio ponderado de los
contratos con destino al mercado regulado. Pero el parágrafo del artículo 25
de la Resolución CREG 174, en la redacción que le dio la Resolución CREG 101 087
de 2025, mantiene el precio de bolsa hasta que la Comisión expida la
metodología de traslado, y los conceptos CREG 3018 y 3023 de abril de 2026
confirman que la transición sigue vigente. En todo el horizonte del trabajo el
excedente sobrante se valora, por tanto, al precio de bolsa de cada hora. Es lo
que hace la autogeneración individual del modelo; el colectivo usaba la media
mensual de la bolsa, que sobrevalora ese excedente porque las horas solares son
las de bolsa más barata, y debe corregirse.

\subsection{Con qué porcentaje reparte el colectivo}
\label{sub:preg-pde}

El artículo 19 de la Resolución CREG 101 072 deja el porcentaje de
distribución de los excedentes a lo que acuerden los integrantes. El modelo
lo calcula en proporción a la generación medida, que en la frontera
principal depende de lo que cada medidor alcanza a ver; con las placas
idénticas, repartir por capacidad instalada daría la misma proporción a las
cinco. Según la respuesta cambian las cuotas de cada institución y, con
ellas, a quién le conviene el colectivo.

\subsection{Si la autogeneración remota se modela como compensación}
\label{sub:preg-agr}

La Resolución CREG 101 099 de 2026 no compensa. La generación se entrega al
mercado mayorista a través de un agente generador, el consumo se atiende
como usuario no regulado con un contrato registrado ante el ASIC, y el
reparto horario con porcentaje fijo sirve solo para devolver el cargo por
confiabilidad. El escenario la modela como una compensación horaria entre
fronteras. Queda por decidir si se reescribe con esos artículos, se conserva
como referencia prospectiva o se retira, teniendo presente que la revisión
externa consideró el régimen cerrado a esta comunidad.

\placeholderfig{Faltan por trasladar aquí las decisiones abiertas que hoy
viven en el registro de correcciones.}


\end{document}
